% Modernised reading — fonds Alexandre Grothendieck, shelfmark 91
% (pages 1-142, the whole folder).
%
% This is an INTERPRETATION, not a transcription. It re-reads the pages in
% today's notation and vocabulary, states the hypotheses the manuscript leaves
% implicit, and drops the critical apparatus so that the mathematics reads
% straight through. Everything the brackets and underlines carried moves into a
% footnote. It is held to being mathematically correct as it stands; where the
% manuscript is loose or mistaken, the true statement is what appears here and
% the footnote says what the page has. In any dispute of substance the
% transcription governs: whoever cites Grothendieck must cite batch-01.fr.tex
% to batch-08.fr.tex (pp. 1-20, 21-40, ..., 121-140, 141-142).
%
% Pass: Opus 5.5 (claude-opus-5-5), 2026-10-17 — first pass, unchecked against the pages by a human.
% The folder was read whole, its eight batches together, from the second-pass
% transcriptions (700 dpi tiles, commit cc339bf); this is one document for the
% shelfmark. It was made on Opus 5.5 and has not been measured against the
% reference reading of folder 115 (made on Opus 5). The literature named in the
% footnotes is cited from memory and was not checked against the sources.
% The transcription still carries many \ill and \uncertain marks, most of them
% in pp. 21-58 and 81-100, where the hand is a fast draft: statements and
% formulas carry, connective prose often does not. Where a section rests on a
% thin reading it says so in the body, not only in a note.
%
% Not transcribed, hence not read: p. 1 and p. 80 are covers (their words are
% used as titles); blanks or unrelated versos 3, 14, 29, 32, 59, 82, 84, 85,
% 87, 89, 91, 93, 95, 97, 99, 101, 102, 111, 116, 118, 120, 122, 124, 126,
% 128, 129, 132, 136, 138, 140, 142 (typed pages of other texts, or nothing).
%
% Spine. One problem, the existence of « spaces of maps », taken three ways:
%   pp. 1-13    subfunctors of the base (monomorphisms): the locus where
%               Y = X, where E is flat, where X -> S is faithfully flat;
%               centralisers, normalisers, transporters; G/N.
%   pp. 15-28   the non-projective method for Hom: the open subfunctor
%               Hom_S(X,Y;phi) where restriction to Z is unramified, its
%               obstruction theory, and the corollaries (Isom, Aut).
%   pp. 30-31   the typed letter of 11.9.1967 to Anantharaman that states it.
%   pp. 33-58   two plans, the formalism of Res, and Weil restriction along a
%               finite locally free T -> S: existence, permanence, ampleness,
%               smoothness, dimension, the relative pi_0.
%   pp. 60-79   « Set up »: a functor W of local nature and its right adjoint
%               W^a; base change (Weil restriction) and Greenberg as cases.
%   pp. 80-121  « Schémas localement compacts »: constructible maps,
%               « ensembles schématiques », the constructible topology.
%   pp. 123-141 « Greenberg-Néron »: thick points, rational maps on Greenberg
%               schemes, the ind-scheme A((t))^*, pro-schemes, jet groups.
%
% Notation fixed for the whole folder. Schemes for « préschémas » (all
% schemes, not separated ones). Res_{T/S}(X) for the Weil restriction of a
% T-scheme X along T -> S, which the pages write prod_{T/S} X/T (pp. 1-28,
% 33), f_*(X/T) (pp. 34-58), phi_{S'/S}(X) (pp. 74-76) and W^a(X) in the case
% of base change (pp. 60-79); Res_{T/S}(g) for the morphism he writes f_*(g).
% Hom_S(X,Y), Isom, Aut underlined: the functors (internal Hom). Hom_S(X,Y;
% phi) as on the page. M_u, M_v the coherent modules of EGA III 7.7.6 attached
% to u, v (pp. 15-20). W for a functor of local nature, W^a its right adjoint;
% Lambda for the ring scheme generating W (the page calls it W too, p. 64).
% « Lisse » for his « simple ». « Pro-constructible » for his « analytique »
% (pp. 103-110), footnoted once. V-check(F) = V(F^vee), the vector bundle
% whose sections are F (p. 37), so that its formula holds.
%
% Substantive departures from the pages, with page numbers:
%  - p. 4: the representing object is a monomorphism of finite type S' -> S;
%    the page's « sous-préschéma (fermé) » for Y closed holds, by EGA IV
%    18.12.6 (proper monomorphism), and is justified here, not on the page.
%  - pp. 11-13: in Cor. 3 (i), f' is a monomorphism of finite presentation,
%    not an immersion as the page has: X = (A^1 - 0) ⊔ {0} -> A^1 satisfies
%    (ii) and is not an immersion. Same correction for the two corollaries of
%    p. 13 (« sous-préschéma » -> monomorphism).
%  - p. 10: he writes F for the coherent sheaf E (and X_. plat for E_. plat).
%  - p. 15, p. 24, letter p. 31: « séparé » for Y (resp. P/X) is supplied, as
%    the letter's inked addition and the « [et séparé ?] » of p. 15 ask.
%  - p. 25: the theorem with Z_1..Z_n is reported as announced; its proof is
%    not on the pages and its sufficiency is not checked here.
%  - p. 48: « f_*(X/T) = Spec(f_*(F)) » read as Res(V-check(F)) =
%    V-check(f_* F), the formula of p. 37.
%  - pp. 50-51: « rang des anneaux locaux artiniens » read as [T_s : k(s)];
%    Cor. 2 reads D(s) = n D'(s), corrected to D'(s) = n D(s) and stated for
%    T_s local (in general D'(s) is a set of sums).
%  - pp. 55-58: f^et_*(T) (relative pi_0) needs nu(s) locally constant (Cor. 3
%    p. 53) to be a finite étale cover; supplied. It is a left adjoint; the
%    name f_* is his and its variance is explained.
%  - p. 69: the page writes « est un W(U') », read W^a(U).
%  - p. 72: the factorisation ends in X_i, read 𝔛_i.
%  - p. 92, Cor. 0: holds for the underlying space, not for the scheme.
%  - p. 109, Cor. 3 (iii): f^{-1}(Y) read f^{-1}(E).
%  - p. 135: Y_j quasi-separated supplied (EGA IV 8.13.1).
%
% Announced and never established in the folder: the two « nous admettrons »
% of pp. 19-20 (compatibility of the obstructions); the question of p. 38
% (does f^* commute with filtered intersections?); the « ??? » of p. 96
% (Cons(lim X_i) = lim Cons(X_i)); the « catégorie géométrique (?) » of p. 98;
% condition (ii) of p. 83 (the sheet is torn); the meaning of « d'indice
% fini » (pp. 130-131); the « Espaces de Néron » of the plan of p. 117; the
% calculation of p. 139, which breaks off; the « ??? » implication of p. 127.
\documentclass[11pt,a4paper]{article}
\input{../preamble/grothendieck}

\folder{91}
\pages{1}{142}
\dating{[à partir de 1964-vers 1970]}
\foldertitle{Autour de Néron / Greenberg-Néron. Foncteurs Hom (méthodes non-projectives) : notes manuscrites (s.d.), lettre (1967) — lecture modernisée du dossier entier}

\begin{document}

\section*{Résumé}

\begin{resume}
Ces feuillets cherchent à construire des \emph{espaces d'applications}.
Étant données deux variétés $X$ et $Y$, on voudrait que l'ensemble de toutes
les applications algébriques de $X$ dans $Y$ soit lui-même une variété, dont
chaque point est une application, et dont les courbes sont des familles
d'applications qui se déforment continûment. Le cas le plus célèbre est celui
des symétries d'une variété : son « groupe d'automorphismes », dont on veut
qu'il soit un groupe algébrique.

Le problème vient de ce que la méthode connue ne marche que pour les variétés
« projectives », celles que l'on peut plonger dans un espace projectif. On y
identifie une application à son graphe, une sous-variété de $X \times Y$, et
l'on dispose d'un espace de toutes les sous-variétés (le schéma de Hilbert).
Sans plongement, il n'y a plus de schéma de Hilbert. Le titre de
l'inventaire le dit : « méthodes non projectives ».

L'idée nouvelle est une \emph{rigidification}. Un polynôme de degré
$d$ est connu dès qu'on connaît ses valeurs en $d+1$ points ; une fonction
holomorphe est connue dès qu'on connaît son développement de Taylor en un
point. De même, une application de $X$ dans $Y$ est, au voisinage d'une
application donnée, déterminée par sa restriction à un petit morceau $Z$ de
$X$ (quelques points épaissis), pourvu que la restriction ne perde aucune
information infinitésimale, ce qu'on dit en demandant qu'elle soit « non
ramifiée ». L'espace des applications de $Z$ dans $Y$ est facile à construire
quand $Z$ est fini ; on en déduit, morceau par morceau, celui des applications
de $X$. C'est ce qu'une lettre de septembre 1967, conservée ici, expose à
un collègue : il y systématise l'étape centrale d'une démonstration de Matsumura
sur le groupe d'automorphismes d'une variété propre.

Pour cela il faut savoir construire l'espace des applications d'un objet
\emph{fini} dans $Y$. C'est la « restriction de Weil ». Dans le cas le plus
familier, une variété définie par des équations à coefficients complexes
devient, en écrivant chaque variable complexe $z = x + iy$, une variété réelle
de dimension double. La moitié du dossier étudie cette opération (quand elle
existe, ce qu'elle préserve) et la généralise : le même mécanisme, appliqué
aux entiers $p$-adiques au lieu des nombres complexes, donne le foncteur de
Greenberg, qui transforme les solutions d'équations modulo $p^n$ en points
d'une variété sur le corps à $p$ éléments. C'est le lien avec le titre,
« Autour de Néron » : la construction des modèles de Néron des variétés
abéliennes passe par là.

Le dernier tiers change d'outil. Pour comparer des points de ces espaces de
Greenberg on veut des applications qui ne sont des morphismes que « par
morceaux » ; il les appelle applications constructibles, et l'espace muni de
la seule topologie qu'elles respectent est un espace compact totalement
discontinu, comme l'ensemble de Cantor. Les feuillets s'arrêtent sur des
calculs de changements de coordonnées, sans avoir rejoint Néron.

Une remarque de la page 135 dit quelque chose de sa manière : là où Serre et
Néron manipulent des systèmes de variétés de plus en plus fines, « il y a
avantage à les interpréter comme étant de véritables schémas » ; autrement dit,
remplacer l'échafaudage par l'objet qu'il approche, et travailler sur l'objet.

Les noms à chercher ensuite : schéma $\mathrm{Hom}$ et théorème de
Matsumura--Oort, critère de représentabilité de Murre, restriction de Weil,
foncteur de Greenberg, modèle de Néron, topologie constructible et parties
pro-constructibles, groupe des lacets de $\mathbf{G}_m$.

\keywords{Hom scheme, Isom scheme, automorphism group scheme, Matsumura-Oort theorem, Murre representability criterion, unramified functor, universal flattening, centralizer and transporter, quotient by a flat normal subgroup, Weil restriction, effective descent, relative connected components, Greenberg functor, Witt vectors, constructible set, constructible topology, pro-constructible set, spectral space, Néron model, ind-scheme, loop group of Gm, pro-scheme, jet group}
\end{resume}

\pagerange{1}{141}
\section*{Le fil du dossier, et les conventions}

Le dossier n'est pas paginé de sa main. Il est fait de suites séparées par
deux feuilles de couverture, « Représentabilité des foncteurs
$\underline{\mathrm{Hom}}$, $\prod_{T/S} X/T$, de Greenberg etc. » (page 1) et
« Schémas localement compacts » (page 80), et par un titre isolé,
« Greenberg -- Néron » (page 123). Beaucoup de versos portent des
dactylographies sans rapport (dont des pages Bourbaki) et ne sont pas
transcrits. L'inventaire date l'ensemble « à partir de 1964 -- vers 1970 » ;
la seule date écrite est celle de la lettre, 11 septembre 1967.

Les stations, dans l'ordre des pages :

\begin{enumerate}
  \item \emph{Pages 1 à 13.} Sous-foncteurs de la base : le lieu où un
  sous-schéma $Y$ de $X$ remplit $X$, le lieu où un faisceau est plat, le lieu
  où $X \to S$ est fidèlement plat. Applications aux centralisateurs,
  normalisateurs, transporteurs, et aux quotients $G/N$.
  \item \emph{Pages 15 à 28.} La méthode non projective pour
  $\underline{\mathrm{Hom}}_S(X,Y)$ : le sous-foncteur ouvert
  $\underline{\mathrm{Hom}}_S(X,Y;\varphi)$, sa théorie d'obstruction, et une
  série de corollaires ($\underline{\mathrm{Isom}}$, $\underline{\mathrm{Aut}}$).
  \item \emph{Pages 30 et 31.} La lettre dactylographiée à Anantharaman qui
  expose cette méthode.
  \item \emph{Pages 33 à 58.} Deux plans d'exposé, le formalisme de la
  restriction de Weil, puis son étude le long d'un morphisme fini localement
  libre.
  \item \emph{Pages 60 à 79.} Un cadre commun : un foncteur $\mathbb{W}$ « de
  nature locale » et son adjoint à droite $\mathbb{W}^{a}$, dont la
  restriction de Weil et le foncteur de Greenberg sont deux cas.
  \item \emph{Pages 80 à 121.} « Schémas localement compacts » : applications
  constructibles, « ensembles schématiques », topologie constructible.
  \item \emph{Pages 123 à 141.} « Greenberg -- Néron » : points épais,
  applications rationnelles entre schémas de Greenberg, l'ind-schéma
  $A \mapsto A((t))^{*}$, pro-schémas, groupes de jets.
\end{enumerate}

Les trois premières stations forment un tout, la quatrième et la cinquième en
forment un autre, et les deux dernières un troisième. Le lien entre eux est
celui que le résumé décrit : la méthode de la deuxième station demande la
restriction de Weil des stations quatre et cinq ; celles-ci, prises dans le
cas des vecteurs de Witt, donnent les schémas de Greenberg des deux
dernières. Le dossier ne referme pas cette boucle : il n'y a ni construction
d'un modèle de Néron, ni retour aux schémas $\underline{\mathrm{Hom}}$ après la
page 58.

\emph{Notations, valables pour tout le dossier.} On écrit « schéma » pour son
« préschéma »\footnote{Le mot « préschéma » est celui d'EGA I (1960) ;
« schéma » y désignait un préschéma séparé. Aucun énoncé ci-dessous ne
suppose la séparation sans le dire.} et « lisse » pour son « simple ». La
\emph{restriction de Weil} d'un $T$-schéma $X$ le long de $f : T \to S$ est le
foncteur sur les $S$-schémas
\[
  \mathrm{Res}_{T/S}(X) : S' \longmapsto \mathrm{Hom}_{T}(T \times_S S', X)
  = \Gamma\bigl(X \times_S S' / T \times_S S'\bigr),
\]
l'ensemble des sections de $X' \to T'$ ; on le note ainsi partout, alors que
les pages l'écrivent successivement $\prod_{T/S} X/T$ (pages 1 à 33),
$f_{*}(X/T)$ (pages 34 à 58), $\varphi_{S'/S}(\mathfrak{X})$ (pages 74 à 76) et
$\mathbb{W}^{a}(\mathfrak{X})$ (pages 60 à 79)\footnote{Les trois premières
écritures sont les siennes ; la quatrième désigne un foncteur plus général,
dont la restriction de Weil est un cas (pages 60 à 79). Le dossier EGA VI
(cote 26) l'écrit, comme ici, $\mathrm{Res}$.}. Pour $g : X \to Y$ un
$T$-morphisme, $\mathrm{Res}_{T/S}(g)$ est le morphisme qu'il note $f_{*}(g)$.
Le symbole souligné $\underline{\mathrm{Hom}}_S(X,Y)$ est le foncteur
$S' \mapsto \mathrm{Hom}_{S'}(X', Y')$, et de même
$\underline{\mathrm{Isom}}$, $\underline{\mathrm{Aut}}$ ; on garde le trait,
qui note aujourd'hui encore le Hom interne. Les formules de la page 34
relient les deux : $\underline{\mathrm{Hom}}_S(X,Y) = \mathrm{Res}_{X/S}(X
\times_S Y)$, et $\mathrm{Res}_{T/S}(X)$ est la fibre de
$\underline{\mathrm{Hom}}_S(T,X) \to \underline{\mathrm{Hom}}_S(T,T)$ au-dessus
de l'identité.

Deux choses portent le même nom dans le dossier et sont tenues distinctes
ici. Le « $f_{*}$ » des pages 34 à 55 est la restriction de Weil, un
adjoint \emph{à droite} ; le « $f_{*}^{\mathrm{\acute{e}t}}$ » des pages 56 à 58
est un adjoint \emph{à gauche}, le schéma des composantes connexes relatives.
Et la lettre $\mathbb{W}$ désigne, page 64, à la fois le foncteur et le schéma en
anneaux qui l'engendre ; on écrit $\Lambda$ pour le second, lettre que les
pages 61 et 65 emploient aussi.

\pagerange{1}{13}
\section*{I. Sous-foncteurs de la base (pages 1 à 13)}

La couverture annonce le programme : « Représentabilité des foncteurs
$\underline{\mathrm{Hom}}$, $\prod_{T/S} X/T$, de Greenberg etc. ». La page 2
n'est qu'une liste de sujets : représentabilité relative de
$\mathrm{Res}_{X/S}(P)$ sans hypothèse projective, « flattening functor »,
passage au quotient, $\underline{\mathrm{Hom}}_S(X,Y)$,
$\underline{\mathrm{Aut}}_S(X)$. Les pages 4 à 13 en traitent la partie la plus
simple : des foncteurs qui sont des \emph{sous-foncteurs} de la base $S$,
c'est-à-dire qui prennent sur chaque $S'$ au plus une valeur.

\subsection*{L'outil : un critère de représentabilité}

Toutes les démonstrations de ces pages et des suivantes vérifient une même
liste de huit conditions, numérotées de 1° à 8°, d'un « théorème de
représentabilité par morphismes non nécessairement séparés » qui n'est pas
énoncé dans le dossier. On reconnaît, par les conditions elles-mêmes et par
la lettre de la page 30 qui l'attribue à « Murre's talk », le critère de
représentabilité des foncteurs non ramifiés exposé par Murre au Séminaire
Bourbaki\footnote{J.~P.~Murre, « Representation of unramified functors.
Applications (according to unpublished results of A.~Grothendieck) »,
Séminaire Bourbaki, exposé 294 (1965). La référence est citée de mémoire. Le
dossier n'en reproduit pas l'énoncé : il en vérifie les conditions une à une
(faisceau pour fpqc ; commutation aux limites inductives d'anneaux ;
commutation aux limites projectives adiques d'anneaux artiniens ; trois
conditions de nature infinitésimale ; deux conditions de « modularité » sur
un schéma local de dimension $1$ muni d'un idéal de carré nul).}. Les conditions
sont, en substance : (1°) $F$ est un faisceau pour la topologie fidèlement
plate quasi-compacte ; (2°) $F$ commute aux limites inductives filtrantes
d'anneaux ; (3°) $F$ commute aux limites projectives adiques d'anneaux locaux
artiniens ; (4° à 6°) des conditions infinitésimales ; (7°) une condition de
recollement : si $T$ est local noethérien, muni d'un idéal $J$ de carré nul,
$T_0 = V(J)$, et $U$ un ouvert schématiquement dense de $T$, le carré
\[
\begin{array}{ccc}
F(T) & \longrightarrow & F(T_0)\\
\downarrow & & \downarrow\\
F(U) & \longrightarrow & F(U_0)
\end{array}
\]
est cartésien ; (8°) si $T$ est intègre et $F(T_0) \neq \emptyset$, $F$ est
représentable au-dessus d'un ouvert dense.

\subsection*{Le lieu où un sous-schéma remplit $X$ (pages 4 à 7)}

\textbf{Théorème.} \emph{Soient $S$ localement noethérien, $f : X \to S$ plat
de type fini, et $Y \to X$ un morphisme de type fini qui est un monomorphisme.
Le sous-foncteur $F$ de $S$, défini par $F(S') \neq \emptyset$ si et seulement
si $Y \times_S S' \to X \times_S S'$ est un isomorphisme, est représentable
(par un monomorphisme de type fini $S_Y \to S$) si et seulement s'il commute
aux limites projectives adiques d'anneaux artiniens. Si $Y \to X$ est une
immersion fermée, $S_Y$ est alors un sous-schéma fermé de
$S$.}\footnote{L'énoncé de la page 4 porte deux additions : « plat », et « de
type fini » pour $Y \to X$ ; une première version, où $Y$ était « un
sous-préschéma \ldots{} de $X$ », est biffée, et une marge ajoute « il suffit
que $Y \to X$ soit un monomorphisme ». La page 5 est une première rédaction
du même énoncé, biffée entière et écrite tête-bêche. Le foncteur est un
sous-foncteur de $S$ ; ce qui le représente est donc un monomorphisme vers
$S$, et c'est ce qu'on écrit. La dernière phrase est sur la page (« alors il
est représentable par un sous-préschéma (fermé) de $S$ ») ; la démonstration
ne la justifie pas. Elle est vraie : quand $Y$ est fermé, la condition 7°
donne, par le raisonnement de la page 6, la partie « existence » du critère
valuatif, et un monomorphisme de type fini propre est une immersion fermée
(EGA IV 18.12.6). Enfin la page énonce le théorème pour un $Y \to X$ de type
fini quelconque, mais sa démonstration ne traite que le cas où le foncteur
est un sous-foncteur de $S$ (« dans le cas d'un monomorphisme $F \to S$ ») ;
on énonce ce qui est démontré.}

Le foncteur est $\mathrm{Res}_{X/S}(Y)$ : une section de $Y' \to X'$, quand
$Y \to X$ est un monomorphisme, existe si et seulement si $Y' = X'$, et elle
est alors unique. Les conditions se vérifient ainsi. La platitude de $X$
donne 4° à 6° : si $A \to A'$ est un homomorphisme injectif d'anneaux
artiniens, $\mathcal{O}_{X_A} \to \mathcal{O}_{X_A} \otimes_A A'$ est injectif,
de sorte que $Y_{A'} = X_{A'}$ entraîne $Y_A = X_A$. Elle donne aussi 7° : si $T$
est l'adhérence schématique de $U$, $X_T$ est l'adhérence schématique de $X_U$
(par platitude), donc un sous-schéma fermé de $X_T$ qui contient $X_U$ est
$X_T$ entier\footnote{La page 6 ajoute deux variantes en interligne : « en
fait, il suffit ici $F(U) \neq \emptyset$ pour $F(T) \neq \emptyset$ », puis
« (6 lorsque $Y$ est supposé fermé dans $X$) ». Le début de la ligne est
illisible.}. La condition 8° vient de la platitude générique (EGA IV 11 : $X_U$
« ensemblistement libre » sur un ouvert dense $U$).

La commutation aux limites adiques s'énonce aussi :

\begin{quote}
(L) pour tout $S$-schéma $T$ spectre d'un anneau local noethérien (qu'on peut
supposer complet, à corps résiduel algébriquement clos), si $Y_T$ contient
un voisinage de la fibre spéciale de $X_T$, alors $Y_T = X_T$.
\end{quote}

\emph{Corollaires.} Si $X$ est propre sur $S$, (L) est automatique. Si $X$ est
un groupe et $Y$ un sous-groupe, $Y_T$ contient un voisinage de la section
unité, donc $Y_T \to X_T$ est un monomorphisme étale, c'est-à-dire une
immersion ouverte ; la vérification de $Y_T = X_T$ se ramène, par
connexité, au cas où $T$ est le spectre d'un anneau de valuation discrète.
D'où la condition, isolée page 8 :

\begin{quote}
(R) pour tout schéma local $T$ sur $S$ et tout sous-groupe ouvert $U$ de
$G_T$ qui contient la fibre spéciale, on a $U = G_T$.
\end{quote}

Pour $G$ de type fini sur $S$ localement noethérien, il suffit de tester (R)
sur les spectres d'anneaux de valuation discrète complets à corps résiduel
algébriquement clos. (R) est vraie quand $G$ est à fibres connexes, ou
quand il existe des sous-schémas $G_i$ finis sur $S$ tels que, pour chaque
$s$, les $G_{i,\bar s} \cdot G^{0}_{\bar s}$ recouvrent $G_{\bar s}$.

\subsection*{Centralisateurs, normalisateurs, transporteurs (pages 8 et 9)}

\textbf{Corollaire.} \emph{Soit $X$ un schéma en groupes plat de type fini sur
$S$ localement noethérien, satisfaisant à (R), et $Y$ un sous-groupe fermé.
Alors $\mathrm{Res}_{X/S}(Y)$ est représentable par un sous-schéma fermé de
$S$.}\footnote{« Plat » manque à l'énoncé de la page 8 ; il est l'hypothèse
du théorème qu'on applique, et la page 9 l'ajoute en interligne pour le
groupe $H$ du corollaire suivant.}

\textbf{Corollaire.} \emph{Soit $u : H \to G$ un homomorphisme de schémas en
groupes sur $S$, avec $G$ localement noethérien, $H$ plat de type fini sur $S$
et satisfaisant à (R) (par exemple à fibres connexes). Alors
$\mathrm{Cent}_G(u)$ est représentable par un sous-schéma de $G$, fermé si $G$
est séparé sur $S$. Si $K$ est un sous-groupe de $G$, le transporteur
$\mathrm{Transp}_G(u, K)$ est représentable, et c'est un sous-schéma (resp.
un sous-schéma fermé) de $G$ si $K$ en est un. Si $K$ est en outre plat et
satisfait à (R), et si $u$ est un monomorphisme, $\mathrm{Transp}_G(H,K)$
est représentable par un schéma de type fini sur $G$ ; en particulier
$\mathrm{Norm}_G(H)$ est représentable, et c'est un sous-schéma (fermé) de $G$
si $H$ l'est.}

On l'obtient en appliquant le théorème au-dessus de $G$ : le centralisateur
est le lieu des $g$ où les deux morphismes $H_G \to G_G$, $u$ et $g u
g^{-1}$, coïncident, c'est-à-dire où le sous-schéma « égaliseur » remplit
$H_G$. Ce sont les énoncés que SGA 3 établit dans les exposés VIII et XI
sous d'autres hypothèses\footnote{Une première version de ce corollaire,
page 8, barrée de deux diagonales, supposait $G$ séparé ; la condition
« ($G$ séparé sur $S$) » migre page 9 dans la conclusion. Le rapprochement
avec SGA 3 est le nôtre ; les pages n'y renvoient pas ici.}.

\subsection*{Le lieu de platitude (pages 10 et 11)}

\textbf{Théorème.} \emph{Soient $f : X \to S$ de type fini, $S$ localement
noethérien, $E$ un $\mathcal{O}_X$-module cohérent. Le sous-foncteur de $S$
\[
  F(S') = \begin{cases} \{\ast\} & \text{si } E_{S'} \text{ est plat sur } S',\\
  \emptyset & \text{sinon,}\end{cases}
\]
est représentable si et seulement s'il commute aux limites projectives
adiques d'anneaux artiniens}\footnote{La page écrit $F_{S'}$ pour le faisceau,
de la même lettre que le foncteur, et, en 4°--7°, « $X_{A'}$ plat » là où il
s'agit de la platitude de $E_{A'}$ ; elle le dit elle-même en 4° (« je veux
dire $E$ plat »). La page 10 s'ouvre sur « Variantes débarrassées
d'hypothèses noethériennes\ldots » et « Variante. Pour un morphisme $Y \to X$
non ramifié », deux titres sans suite.}. \emph{La condition s'écrit : pour tout
$T$ spectre d'un anneau local noethérien, si $E_T$ est plat sur $T$ au
voisinage de la fibre générique, il est plat. Elle est satisfaite si $f$ est
propre, et si $E = \mathcal{O}_X$ avec $X$ un schéma en groupes satisfaisant à
(R).}

C'est ce que la page 2 appelait « flattening functor ». Le cas propre est
aujourd'hui le théorème de l'aplatissement universel, dû à Mumford dans le
cas projectif et à Raynaud et Gruson en général\footnote{Le rapprochement est
le nôtre et la référence citée de mémoire : M.~Raynaud, L.~Gruson, « Critères
de platitude et de projectivité », \emph{Invent. Math.} 13 (1971). Une marge
de la page 11 ajoute qu'on peut supposer $\dim T = 1$, $T$ intègre, donc le
spectre d'un anneau de valuation discrète ; deux mots y sont illisibles.}.

\subsection*{Le lieu où un morphisme est fidèlement plat, et les quotients (pages 11 à 13)}

\textbf{Corollaire 3.} \emph{Soit $f : X \to S$ de présentation finie.
Conditions équivalentes :}
\begin{enumerate}
  \item[(i)] \emph{$f$ se factorise en $X \xrightarrow{f''} S' \xrightarrow{f'} S$,
  avec $f''$ fidèlement plat et $f'$ un monomorphisme ;}
  \item[(i bis)] \emph{de même avec $f''$ seulement plat ;}
  \item[(ii)] \emph{la projection $\mathrm{pr}_1 : X \times_S X \to X$ est plate.}
\end{enumerate}
\emph{Sous ces conditions la factorisation (i) est unique à isomorphisme unique
près, $S' = X/R$ avec $R = X \times_S X$, et $S'$ est de présentation finie sur
$S$.}\footnote{La page demande que $f'$ soit une \emph{immersion} (« $f'$
\emph{imm.} », souligné). C'est trop : le morphisme $X = (\mathbf{A}^1 - \{0\})
\sqcup \{0\} \to \mathbf{A}^1$ satisfait à (ii) (au-dessus de chaque morceau,
$X \times_S X \to X$ est un isomorphisme), et c'est un monomorphisme
bijectif qui n'est pas une immersion ; ici $S' = X$. La démonstration de la
page 12 ne fournit qu'un monomorphisme : le foncteur qu'elle
représente est le sous-foncteur de $S$ des $T$ tels que $X_T \to T$ soit
couvrant pour fpqc, et son représentant est donc un monomorphisme vers $S$.}

La démonstration est celle de la page 12 : (i) $\Rightarrow$ (i bis)
$\Rightarrow$ (ii) est immédiat, puisque $X \times_S X = X \times_{S'} X$ quand
$S' \to S$ est un monomorphisme ; l'unicité vient de ce que $X \to S'$ est un
épimorphisme effectif pour fpqc. Pour (ii) $\Rightarrow$ (i), on se ramène à
$S$ affine noethérien et l'on applique le théorème au foncteur « $X_T \to T$
fidèlement plat » ; sous (ii), la condition (L) est satisfaite, car un point
$x$ de la fibre spéciale donne un anneau local $\mathcal{O}_{X,x}$ fidèlement
plat sur celui de $T$.

\textbf{Corollaires} (page 13). \emph{(a) Soient $G$ un schéma en groupes de
présentation finie sur $S$ et $P$ un $S$-schéma de présentation finie sur
lequel $G$ opère de façon que $P$ soit formellement principal homogène. Il
existe un monomorphisme $S' \to S$ par lequel $P \to S$ se factorise et au-dessus
duquel $P$ est un torseur (fidèlement plat) sous $G$ ; $S'$ est unique à
isomorphisme unique près. (b) Soit $u : G \to H$ un homomorphisme de schémas
en groupes de présentation finie, dont le noyau $N$ est plat (et de présentation
finie) sur $S$. Alors le quotient $G/N$ existe, comme schéma, et $G \to G/N
\to H$ factorise $u$ en un épimorphisme fidèlement plat suivi d'un
monomorphisme.}

Pour (b) : $G \times_H G \simeq G \times_S N$ est plat sur $G$, et le corollaire 3
s'applique à $G \to H$\footnote{La page dit seulement « $G/N$ existe ». La
factorisation par un monomorphisme est ce que le corollaire 3, corrigé,
fournit ; avec l'énoncé de la page on aurait conclu à tort que $G/N \to H$
est une immersion. Le (a) dit sur la page « il existe un sous-préschéma
$S' \hookrightarrow S$ » : même correction.}.

\pagerange{15}{24}
\section*{II. Le foncteur $\underline{\mathrm{Hom}}_S(X, Y; \varphi)$ (pages 15 à 24)}

\subsection*{La donnée, et le sous-foncteur ouvert (pages 15 et 16)}

On se donne sur $S$ un morphisme $\varphi : Z \to X$ entre deux $S$-schémas
$X$ et $Z$ plats et propres, et un $S$-schéma $Y$ de type fini et
séparé\footnote{Page 15 : « $f$, $h$ plat et propre sur $S$ loc.\ noeth. »,
« $g$ de type fini [et séparé ?] » ; la séparation est restée une question.
La lettre des pages 30 et 31 l'ajoute à l'encre (« and separated ») à
l'énoncé correspondant, et on la suppose ici. Une marge ajoute « suffit si
$Z \to X$ est surjectif », sans dire à quoi ; un N.B. dit que l'hypothèse
« $S$ localement noethérien » peut être remplacée par des hypothèses de
présentation finie sur $X$, $Y$, $Z$.}. La composition avec $\varphi$ définit
\[
  \varphi^{\sharp} : \underline{\mathrm{Hom}}_S(X,Y) \longrightarrow
  \underline{\mathrm{Hom}}_S(Z,Y), \qquad u \longmapsto u\varphi .
\]
Pour $u : X' \to Y'$ au-dessus de $S'$, posons $v = u\varphi'$. Comme $X'$ et
$Z'$ sont propres et plats, il existe des $\mathcal{O}_{S'}$-modules cohérents
$M_u$, $M_v$, dont la formation commute au changement de base, tels que pour
tout $\mathcal{O}_{S'}$-module quasi-cohérent $J$
\[
  \mathrm{Hom}_{\mathcal{O}_{X'}}\bigl(u^{*}\Omega^1_{Y'/S'},\, \mathcal{O}_{X'}
  \otimes J\bigr) \simeq \mathrm{Hom}_{\mathcal{O}_{S'}}(M_u, J),
\]
et de même pour $v$ et $Z'$ (EGA III 7.7.6). L'application naturelle entre les
deux groupes de gauche provient d'un homomorphisme $M_v \to M_u$. Le
morphisme $\varphi^{\sharp}$ est « non ramifié en $u$ » au-dessus d'un point
$s$ si l'application tangente
$\mathrm{Hom}(u_s^{*}\Omega^1, k(s)) \to \mathrm{Hom}(v_s^{*}\Omega^1, k(s))$
est injective, c'est-à-dire si $M_v \otimes k(s) \to M_u \otimes k(s)$ est
surjectif ; par Nakayama, la condition est ouverte. On note
$\underline{\mathrm{Hom}}_S(X,Y;\varphi)$ le sous-foncteur ouvert des $u$ pour
lesquels $M_v \to M_u$ est surjectif.

\textbf{Théorème} (page 16). \emph{Le morphisme
$\underline{\mathrm{Hom}}_S(X,Y;\varphi) \to \underline{\mathrm{Hom}}_S(Z,Y)$ est
représentable par des morphismes non ramifiés. En particulier, si
$\underline{\mathrm{Hom}}_S(Z,Y)$ est représentable, $\underline{\mathrm{Hom}}_S(X,Y;
\varphi)$ l'est aussi.}

\subsection*{La démonstration : l'obstruction (pages 16 à 20)}

On fixe $v : Z \to Y$ et l'on représente le foncteur des $u$ tels que
$u\varphi = v$ et que $M_v \to M_u$ soit surjectif, par le critère des
pages 4 à 7 (version non ramifiée et séparée). Toutes les conditions sont
« habituelles » sauf les deux de modularité, et la seule qui demande un
calcul est 7° : $S = \mathrm{Spec}\,A$ local noethérien de dimension $1$
(dans le calcul, l'anneau d'une uniformisante $\pi$), $J$ un idéal de carré
nul, $S_0 = V(J)$, $\eta$ le point générique ; on a $u_0 : X_0 \to Y_0$ avec
$u_0\varphi_0 = v_0$, et un prolongement $u_\eta$ avec $u_\eta \varphi_\eta = v_\eta$.
Il s'agit de prolonger $u_0$ en $u$ avec $u\varphi = v$.

Les prolongements de $u_0$ (resp. $v_0$) sur $S$ et sur le localisé en $\eta$
forment un carré d'ensembles principaux homogènes sous le carré de groupes
\[
\begin{array}{ccc}
\mathrm{Hom}(M_{u_0}, J) & \longrightarrow & \mathrm{Hom}(M_{u_0}, J_\eta)\\
\downarrow & & \downarrow\\
\mathrm{Hom}(M_{v_0}, J) & \longrightarrow & \mathrm{Hom}(M_{v_0}, J_\eta)
\end{array}
\]
et l'on cherche un élément du coin supérieur gauche relevant les deux
éléments donnés aux coins voisins\footnote{L'astérisque de « $M_{u_0}^{*}$ »
au coin supérieur gauche est sur la page ; il ne correspond à rien et on
l'omet. L'indice que la main forme comme un $y$ est lu $\eta$, le point
générique nommé plus haut ; dans les pages 19 et 20, $J_n$ et $J_\eta$ se
ressemblent et sont départagés par le sens.}. Il suffit que l'ensemble des
prolongements de $u_0$ sur $S$ soit non vide : le carré d'ensembles est alors
isomorphe au carré de groupes, et celui-ci est cartésien parce que
$M_{v_0} \to M_{u_0}$ est surjectif et que $J$ est sans torsion sur $A$, ce qui
rend $\mathrm{Hom}(N, J) \to \mathrm{Hom}(N, J_\eta)$ injectif\footnote{La page dit
« (2) est cartésien, grâce au fait que $M_{v_0} \to M_{u_0}$ est surjectif » ;
l'injectivité de la localisation, qu'on ajoute, est ce qui fait marcher
l'argument, et elle tient dans le cas où $A$ est intègre, celui du calcul.}.

L'obstruction à l'existence de $u$ est une classe
$\xi \in \mathrm{Ext}^1_{\mathcal{O}_{P_0}}(I_{\Gamma_0}, \mathcal{O}_{\Gamma_0}
\otimes J)$, où $P_0 = X_0 \times_{S_0} Y_0$, $\Gamma_0$ est le graphe de $u_0$ et
$I_{\Gamma_0}$ son idéal. Elle s'annule en $\eta$, donc $\pi^{n+1}\xi = 0$ pour un
$n$. La suite exacte $0 \to J \xrightarrow{\pi^{n+1}} J \to J_n \to 0$, avec
$J_n = J/\pi^{n+1}J$, permet de la relever en une classe de
$\mathrm{Hom}(I_{\Gamma_0}/I_{\Gamma_0}^2, \mathcal{O}_{\Gamma_0} \otimes J_n) =
\mathrm{Hom}(M_{u_0}, J_n)$ modulo l'image de $\mathrm{Hom}(M_{u_0}, J)$, d'où une
classe $\xi'' \in \mathrm{Ext}^1_{A_0}(M_{u_0}, J)$, tuée par une puissance de
$\pi$, dont l'annulation équivaut au prolongement de $u_0$. De même pour $v_0$
une classe $\eta''$, qui est nulle puisque $v$ existe. Avec la suite exacte
$0 \to N \to M_{v_0} \to M_{u_0} \to 0$, $\xi''$ provient d'une classe
$\xi''' \in \mathrm{Hom}(N, J)/\mathrm{Im}\,\mathrm{Hom}(M_{v_0}, J)$ ; la compatibilité
des prolongements génériques donne une classe $\eta'$ dans
$\mathrm{Hom}(N, J_\eta)/\mathrm{Im}\,\mathrm{Hom}(M_{v_0}, J)$, nulle par hypothèse, et
si $\eta' = \xi'''$ on conclut.

Deux pas sont admis, et la page le dit : que $\eta''$ soit l'image de
$\xi''$ par $M_{v_0} \to M_{u_0}$, et que $\eta' = \xi'''$. La démonstration
n'est donc pas achevée sur ces feuillets\footnote{« Nous admettrons » deux
fois, pages 19 et 20. « $\eta''$ est nul » est souligné.}.

\subsection*{Une variante sous hypothèse de rigidité (pages 21 à 24)}

La page 21 se propose de « simplifier les raisonnements, en éliminant les
hypothèses de propreté sur $X$, $Z$ et le résultat délicat de l'existence de
$M_u$, $M_v$ », en imposant une hypothèse de rigidité : l'application
\[
  \mathrm{Hom}\bigl(u_0^{*}\Omega^1_{Y_0/S_0}, \mathcal{O}_{X_0} \otimes G\bigr)
  \longrightarrow
  \mathrm{Hom}\bigl(v_0^{*}\Omega^1_{Y_0/S_0}, \mathcal{O}_{Z_0} \otimes G\bigr)
\]
injective pour tout $G$ quasi-cohérent sur $S_0$\footnote{La page écrit
$\omega^{\circ}$ pour les deux images inverses et $\Omega^1_{X_0/S_0}$ au second
membre ; on lit $v_0^{*}\Omega^1_{Y_0/S_0}$, seul sens que la flèche puisse
avoir. Une marge porte, au crayon et entourée, « SGA 3 III 4.2 », le
paragraphe sur les relèvements infinitésimaux de morphismes. Le reste de la
page est biffé.}. Les pages 22 à 24 refont alors le calcul : les relèvements
de $u_0$ correspondent aux scindages d'une extension
$0 \to \mathcal{O}_{\Gamma_0} \otimes J \to E \to I_{\Gamma_0} \to 0$ sur $P_0$,
ceux de $v_0$ à ceux de l'extension qui s'en déduit sur $Z_0 \times_{S_0} Y_0$,
laquelle est déjà scindée. Les foncteurs $\mathrm{Ext}^1$ en $G$ sont
représentés, d'où un diagramme de modules $L \to M \to N \to 0$ au-dessus de
$0 \to L' \to M' \to N' \to 0$, et la condition d'existence d'un relèvement
compatible, pour un $\lambda : J \to G$ donné, est l'annulation de $\lambda$ sur
un sous-module $R = R_1 + R_2$ de $J$ ; autrement dit, le foncteur en $G$ est
représenté par $J/R$, le conoyau d'une flèche
$\rho : L \oplus M' \to J$, et sa formation commute aux extensions plates de
la base. C'est ce qui donne les conditions 7° et 8°. La surjectivité de
$(i,\beta) : L \oplus M' \to M$, notée en marge, est la condition de non
ramification\footnote{La lecture de ces trois pages est mince : la prose qui
relie les formules est en grande partie illisible ou incertaine, et l'on
ne reconstruit que la structure du calcul : deux extensions, leurs
$\mathrm{Ext}^1$ représentés, un quotient $J/R$. Les flèches $\alpha$, $\beta$ du
diagramme (D) portent une pointe à chaque bout sur la page.}.

\pagerange{24}{28}
\section*{III. Les théorèmes de représentabilité (pages 24 à 28)}

\textbf{Corollaire} (page 24). \emph{Soit $k$ un corps algébriquement clos,
$X$ propre sur $k$ et $P \to X$ localement de type fini et séparé. Alors
$\mathrm{Res}_{X/k}(P)$ est représentable par un schéma localement de type fini.
Si $P$ est quasi-projectif sur $X$, ou seulement à fibres quasi-projectives,
toute partie ouverte quasi-compacte en est quasi-projective sur $k$ ; il suffit
que pour tout $x \in \mathrm{Ass}\,\mathcal{O}_X$ il existe $y \in \overline{\{x\}}$
tel que la fibre $P_y$ soit quasi-projective.}\footnote{En marge :
« (déjà difficile !) ». « Séparé » est supposé ici, comme pour $Y$ à la page 15
et dans la lettre ; voir la note correspondante.}

\textbf{Théorème} (page 25)\textbf{,} \emph{tel qu'il est annoncé.} Soient $X$
propre, plat et de présentation finie sur $S$, $P \to X$ localement de
présentation finie, et $Z_1, \ldots, Z_n$ des sous-schémas fermés de $X$, plats
et de présentation finie sur $S$. On suppose :
\begin{enumerate}
  \item[(a)] chaque $\mathrm{Res}_{Z_i/S}(P_{Z_i})$ est représentable (par
  exemple $Z_i$ radiciel sur $S$, ou $Z_i/S$ et $P_{Z_i}/Z_i$ projectifs) ;
  \item[(b)] $X$ est normalement plat le long de chaque $Z_i$ (par exemple
  $Z_i$ lisse sur $S$ et $X$ lisse sur $S$ aux points de $Z_i$) ;
  \item[(c)] pour tout $s \in S$ et tout point associé $x$ de $X_{\bar s}$,
  $\bigcup_i Z_{i,\bar s}$ rencontre l'adhérence de $x$.
\end{enumerate}
Alors (i) $\mathrm{Res}_{X/S}(P)$ est représentable par un schéma localement de
présentation finie sur $S$ ; (ii) pour toute partie ouverte quasi-compacte $U$
de celui-ci, $U \to \prod_i \mathrm{Res}_{Z_i/S}(P_{Z_i})$ est quasi-affine ;
(iii) si les $\mathrm{Res}_{Z_i/S}(P_{Z_i})$ sont quasi-projectifs sur $S$ (par
exemple $Z_i$ fini et $P_{Z_i}$ projectif), toute partie quasi-compacte de
$\mathrm{Res}_{X/S}(P)$ est quasi-projective sur $S$, et il existe un fibré en
droites sur $\mathrm{Res}_{X/S}(P)$ ample relativement à $S$ sur toute partie
quasi-compacte.

La démonstration n'est pas dans le dossier, et nous ne vérifions pas que
les hypothèses (a) à (c) suffisent\footnote{Le mécanisme est celui de la
page 15 et de la lettre : recouvrer $\underline{\mathrm{Hom}}$ par les ouverts
$\underline{\mathrm{Hom}}(X,Y;\varphi)$. Avec une famille \emph{fixe} de $Z_i$, la
restriction des sections de $u^{*}T_Y$ à $\bigcup Z_i$ n'a pas de raison
d'être injective ; il faut vraisemblablement faire jouer aussi les voisinages
infinitésimaux des $Z_i$, et c'est le rôle probable de (b), puisque la normale
platitude rend ces voisinages plats sur $S$. C'est une conjecture de lecture.
Au (i), « séparé » est biffé dans la conclusion ; au (ii), le mot « ouverte »
et la parenthèse « (qui est quasi-séparé) » sont des lectures incertaines.}. Il
est suivi, pages 26 à 28, d'une liste de corollaires :

\begin{itemize}
  \item \emph{Cor. 1} est contenu dans le théorème.
  \item \emph{Cor. 2.} $k$ un corps quelconque, $X$ propre sur $k$, $P/X$
  localement de type fini : si tout point associé $x_i$ de $X$ a dans son
  adhérence un point $x'_i$ à corps résiduel radiciel sur $k$ (plus
  généralement, tel que $\mathrm{Res}_{k(x'_i)/k}(P_{x'_i})$ soit représentable),
  $\mathrm{Res}_{X/k}(P)$ est représentable ; c'est le cas si $X$ est irréductible
  et a un point rationnel.
  \item \emph{Cor. 3.} Sans hypothèse sur les points, la représentabilité a
  lieu après une extension finie séparable $k'/k$.
  \item \emph{Cor. 4.} $P/X$ à fibres quasi-projectives, $X/k$ propre :
  $\mathrm{Res}_{X/k}(P)$ existe et toute partie ouverte quasi-compacte est
  quasi-projective.
  \item \emph{Cor. 5.} $X$ propre et plat sur $S$ noethérien intègre : il
  existe un ouvert $U \neq \emptyset$ et un revêtement fini étale surjectif $U'
  \to U$ au-dessus duquel $\mathrm{Res}$ est représentable.
  \item \emph{Cor. 6.} $X \to S$ propre et plat à fibres géométriques
  intègres, muni d'une section le long de laquelle il est lisse :
  $\mathrm{Res}_{X/S}(P)$ est représentable.
  \item \emph{Cor. 7.} $X \to S$ lisse et propre : représentable localement
  pour la topologie étale ; pour $X$ seulement propre, plat, à fibres sans
  cycles associés immergés, localement pour la topologie fidèlement plate
  quasi-finie.
  \item \emph{Cor. 8.} $X$ et $Y$ propres sur un corps $k$ :
  $\underline{\mathrm{Isom}}_k(X,Y)$ est représentable ; en particulier
  $\underline{\mathrm{Aut}}_k(X)$.
  \item \emph{Cor. 9.} $X$ propre et plat sur $S$ noethérien intègre : il
  existe $U \neq \emptyset$ et $U' \to U$ fini étale surjectif tels que
  $\underline{\mathrm{Aut}}_S(X) \times_S U'$ soit représentable. « Il est
  plausible qu'on puisse même prendre $U' = U$, c'est lié à une question
  d'effectivité de données de descente » ; en marge, « cf.\ Raynaud ».
\end{itemize}

Le corollaire 8 est le théorème de Matsumura et Oort : pour $X$ propre sur un
corps, $\underline{\mathrm{Aut}}_k(X)$ est un schéma en groupes localement de type
fini\footnote{H.~Matsumura, F.~Oort, « Representability of group functors, and
automorphisms of algebraic schemes », \emph{Invent. Math.} 4 (1967), cité de
mémoire. La lettre de la page 30 attribue le théorème à Matsumura seul.
Aujourd'hui, pour $X$ propre et plat de présentation finie sur $S$ et $Y$
séparé, $\underline{\mathrm{Hom}}_S(X,Y)$ est en général un espace algébrique
(Artin) ; ces feuillets cherchent un \emph{schéma}, et c'est
ce qui oblige aux hypothèses sur les points et aux extensions de la base.}.
Le dossier EGA VI (cote 26) annonce au § 22 des « théorèmes du type
Matsumura » pour $\mathrm{Res}_{Y/S}(X)$ et les foncteurs voisins ; ces pages
sont, en partie, ce que ce paragraphe aurait contenu, par une autre méthode
que le schéma de Hilbert que la lecture du dossier 26 lui prête.

\pagerange{30}{31}
\section*{IV. La lettre à Anantharaman (pages 30 et 31)}

Deux feuillets dactylographiés en anglais, sans signature, datés à la main
« 11.9.1967 », adressés « Dear Anantharaman ». Ils sont de Grothendieck : le
« je » y expose la méthode des pages 15 à 28 et propose d'envoyer « a
photocopy of the statement of the general theorem », c'est-à-dire du
théorème de la page 25\footnote{L'attribution se tire du contenu, la lettre
n'étant pas signée ; c'est vraisemblablement un double conservé par
l'expéditeur. La transcription la reproduit telle qu'elle est tapée, fautes
de frappe comprises (« standart », « descebt », « resultst »), avec deux
insertions dactylographiées en interligne et une à l'encre.}.

En substance, Matsumura a prouvé que, pour $X$ propre sur un corps
$k$, $\underline{\mathrm{Aut}}_{X/k}$ est représentable par un schéma en groupes
localement de type fini ; la lettre propose d'en systématiser l'étape centrale.
Pour $\varphi : Z \to X$ entre $S$-schémas propres, plats et de présentation
finie, et $Y$ localement de présentation finie et séparé, on a le
sous-foncteur ouvert $\underline{\mathrm{Hom}}_S(X,Y;\varphi)$ où $\varphi^{\sharp}$
est non ramifié, et le résultat de l'exposé de Murre montre que
$\underline{\mathrm{Hom}}_S(X,Y;\varphi) \to \underline{\mathrm{Hom}}_S(Z,Y)$ est
représentable par des morphismes non ramifiés, séparés, localement de
présentation finie. Pour représenter $\underline{\mathrm{Hom}}_S(X,Y)$, on cherche
des $\varphi_i : Z_i \to X$ dont les ouverts $\underline{\mathrm{Hom}}_S(X,Y;\varphi_i)$
recouvrent $\underline{\mathrm{Hom}}_S(X,Y)$ comme faisceau fpqc et dont les
$\underline{\mathrm{Hom}}_S(Z_i, Y)$ sont représentables. Sur un corps
algébriquement clos on prend pour $Z_i$ tous les sous-schémas finis de $X$ à
points radiciels ; sur un corps quelconque, cela marche après extension finie,
d'où le théorème de Matsumura par descente.

Elle ajoute trois remarques. $\underline{\mathrm{Hom}}_S(X,Y)$ n'est pas toujours
représentable, même quand $X$ est le spectre d'une extension quadratique de
$k$ et $Y$ propre non projectif ; c'est la non-existence de certaines
restrictions de Weil, dont les pages 34 et 35 donnent un exemple. Sur une base
quelconque, les points radiciels sont remplacés par des sous-schémas plats :
$\underline{\mathrm{Hom}}_S(X,Y)$ est représentable si $X$ a des fibres
géométriques intègres et une section le long de laquelle il est lisse, et
localement pour la topologie étale si $X$ a des fibres géométriques
réduites\footnote{C'est le corollaire 7 de la page 27, qui demande « $X \to S$
simple propre » ; la lettre dit « reduced geometric fibers ». Les deux ne
coïncident pas, et le dossier ne contient pas l'énoncé intermédiaire.} ; et il
est représentable si $Y$ est quasi-projectif sur un corps.

\pagerange{33}{36}
\section*{V. Deux plans, le formalisme, et le contre-exemple de Nagata (pages 33 à 36)}

La page 33 est le plan d'un exposé « Schémas Hom etc. » : (1) formalisme des
$\underline{\mathrm{Hom}}$ et des restrictions de Weil ; (2) foncteurs de
Serre--Lang et de Greenberg ; (3)--(4) existence ; (5) application aux
immersions et morphismes de type fini ; (6) morphismes quasi-projectifs ;
(7) morphismes lisses et étales ; [8] morphismes quasi-étales ; (9) cas d'un
$S' \to S$ étale, « à reporter à la fin » ; [10] compléments dans le cas de
Greenberg\footnote{Plusieurs mots du plan sont illisibles (rubriques 3 et 5). Le
$\underline{W}^{\circ}(X)$ des rubriques 2 et 3 est lisible, son sens non.}. C'est,
à peu près, l'ordre des pages 37 à 79.

La page 34 pose le formalisme. Pour $T \to S$, $X$ sur $T$, $Y$ sur $S$ :
\begin{align*}
  \mathrm{Res}_{T/S}(X) &= \underline{\mathrm{Hom}}_S(T, X) \times_{\underline{\mathrm{Hom}}_S(T,T)} S,\\
  \underline{\mathrm{Hom}}_{T/S}(X, Y_T) &= \mathrm{Res}_{T/S}\bigl(\underline{\mathrm{Hom}}_T(X, Y_T)\bigr),\\
  \underline{\mathrm{Hom}}_S(X,Y) &= \mathrm{Res}_{X/S}(X \times_S Y).
\end{align*}
Et une formule de descente : si $X$ sur $T$ est muni d'une donnée de descente
relativement à $T \to S$ et si les restrictions de Weil existent, l'objet
descendu existe et
\[
  X_0 = \mathrm{Res}_{T/S}(X) \times_{\mathrm{Res}_{T \times_S T/S}(X \times_S T)} \mathrm{Res}_{T/S}(X),
\]
les deux flèches étant l'image inverse par $\mathrm{pr}_1$ et l'image inverse
par $\mathrm{pr}_2$ suivie de la donnée de descente\footnote{À première vue
on attendrait un égalisateur et non un produit fibré. Les deux coïncident
ici : un couple $(s, t)$ de sections avec $\mathrm{pr}_1^{*}s =
\theta(\mathrm{pr}_2^{*}t)$ donne, en restreignant à la diagonale, $s = t$, la
donnée de descente $\theta$ étant l'identité sur la diagonale.}.

\emph{Le contre-exemple} (pages 34 et 35). Soit $X$ un schéma lisse sur un corps
algébriquement clos $k$, muni d'une action libre de $G = \mathbf{Z}/2$ sans
quotient schématique, parce qu'une orbite n'est contenue dans aucun ouvert
affine\footnote{La page dit « contre-exemple de Nagata ». L'exemple d'une
involution libre sur une variété lisse propre sans quotient schématique est
souvent attribué à Hironaka ; nous gardons le nom de la page. Le raisonnement
de la page 35 est presque entièrement illisible entre les formules ; on n'en
garde que ce qu'il affirme : une orbite fermée contenue dans aucun ouvert
affine.}. Soient $L$ une extension de $k$ sur laquelle $G$ opère fidèlement et
$K = L^{G}$. L'action diagonale sur $X \otimes_k L$ est une donnée de descente
galoisienne relativement à $L/K$, non effective. Par la formule précédente :

\textbf{Corollaire.} \emph{$\mathrm{Res}_{L/K}(X_L)$ n'est pas représentable ; a
fortiori $\underline{\mathrm{Hom}}_K(\mathrm{Spec}\,L, X_L)$ ne l'est pas.}

C'est l'exemple qu'annonçait la lettre\footnote{Cf. S.~Bosch, W.~Lütkebohmert,
M.~Raynaud, \emph{Néron Models} (1990), § 7.6, pour la condition d'existence
des restrictions de Weil et ce type de contre-exemple ; cité de mémoire.}.

La page 36 est un second plan, en paragraphes : § 1 descente d'objets plats
par morphismes quelconques ; § 2 descente d'objets quelconques par morphismes
plats ; § 3 application aux morphismes non ramifiés ; § 4 cas élémentaires
(quotients $X/G$, descente galoisienne, le contre-exemple de Nagata,
descente plate finie par opérateurs différentiels, cas de Cartier) ; § 5
$\underline{\mathrm{Hom}}_S(X,Y)$ et restriction de Weil, existence pour $X$ fini et
plat, les $\check{\mathbf{V}}(\mathcal{F})$ ; § 6 applications (géométrie
différentielle, Picard, Greenberg) ; § 7 hypothèses sur $S$, critère de
représentabilité ; § 8 théorème fondamental d'existence : théorie formelle des
modules, restriction de Weil pour $X$ projectif sur un $Y$ plat et propre,
Picard. C'est le plan d'EGA VI dans une de ses versions, et il recoupe le plan
du dossier 26 (descente, puis représentabilités)\footnote{Les §§ 9 et 10 sont
sans intitulé lisible ; une accolade réunit les §§ 8 à 10 avec en marge
« ensemble ? ».}.

\pagerange{37}{58}
\section*{VI. La restriction de Weil le long d'un morphisme fini (pages 37 à 58)}

Dans toute cette partie $f : T \to S$ est fini et localement libre, de rang $n$
quand il est constant.

\subsection*{Fibrés vectoriels, et une question (pages 37 et 38)}

\textbf{Proposition.} \emph{Pour $\mathcal{F}$ localement libre de type fini sur
$T$, on a un isomorphisme fonctoriel, compatible au changement de base,
\[
  \mathrm{Res}_{T/S}\bigl(\check{\mathbf{V}}(\mathcal{F})\bigr) =
  \check{\mathbf{V}}\bigl(f_{*}\mathcal{F}\bigr),
\]
où $\check{\mathbf{V}}(\mathcal{F}) = \mathbf{V}(\mathcal{F}^{\vee})$ est le fibré
vectoriel dont le faisceau des sections est $\mathcal{F}$.}\footnote{La notation
$\check{\mathbf{V}}$ est la sienne, sans définition dans le dossier. Avec le
$\mathbf{V}(\mathcal{F}) = \mathrm{Spec}\,\mathrm{Sym}\,\mathcal{F}$ d'EGA II, dont les
sections sont les formes linéaires sur $\mathcal{F}$, la formule serait fausse
($\mathrm{Hom}_T(\mathcal{F}, \mathcal{O}_T)$ n'est pas
$\mathrm{Hom}_S(f_{*}\mathcal{F}, \mathcal{O}_S)$) ; elle est vraie, et immédiate,
pour le fibré des sections, et c'est le sens qu'on donne ici au signe.}

La page 38 pose une question qui reste ouverte sur la feuille. Soient $f : X
\to Y$ plat et propre, $Y$ local noethérien, $F$ cohérent sur $Y$ et $F'_1$ un
sous-faisceau cohérent de $f^{*}F$ : y a-t-il un plus petit sous-faisceau
cohérent $F_1$ de $F$ avec $f^{*}F_1 \supset F'_1$ ? Par platitude, la famille
des $F_i$ convenables est stable par intersection finie ; tout revient à
savoir si $f^{*}$ commute aux intersections filtrantes décroissantes, et la
page écrit : « est-il vrai que $\bigcap f^{*}(F_i) = 0$ ? »\footnote{La
démonstration esquissée en bas de page est presque illisible. La propriété
demandée est ce qu'on appelle aujourd'hui une platitude « par intersections »
(Raynaud--Gruson) ; nous ne savons pas si elle tient sous ces seules
hypothèses et ne tranchons pas.}

\subsection*{Existence (pages 39 à 42)}

\textbf{Proposition 2 bis.} \emph{Soit $X = \bigcup_i U_i$ un recouvrement
ouvert tel que, pour tout $s \in S$, toute section de $X_s \to T_s$ définie
sur une extension finie de $k(s)$ ait son image dans un des $U_i$. Si les
$\mathrm{Res}_{T/S}(U_i)$ existent, $\mathrm{Res}_{T/S}(X)$ existe et en est la
réunion, les $\mathrm{Res}_{T/S}(U_i)$ y étant ouverts.}

La démonstration (pages 40 et 41) recolle les $\mathrm{Res}_{T/S}(U_i)$ le long de
leurs intersections en un schéma $H$, puis montre que $H \times_S T \to X$ est
universel : une section $\sigma$ de $X_{S'} \to T_{S'}$ envoie, au voisinage de
chaque point $s'$, $T_{s'}$ dans un $U_i$, donc tout un voisinage, parce que
$T$ est fini, donc fermé\footnote{La page 39 donne d'abord une
« Proposition 2 », remplacée en marge par la 2 bis ; la démonstration des
pages 40 et 41 est très surchargée et sa prose en grande partie incertaine.
On n'en garde que l'argument que les formules portent.}.

\textbf{Théorème} (page 41). \emph{Si, pour tout $s \in S$, toute multisection
de $X_s$ sur $T_s$ définie sur une extension finie de $k(s)$ est contenue dans
un ouvert affine de $X$, alors $\mathrm{Res}_{T/S}(X)$ existe et est la réunion
des $\mathrm{Res}_{T/S}(U)$, $U$ parcourant les ouverts affines de $X$.}

Le cas affine est immédiat : la restriction de Weil d'un $T$-schéma affine est
un $S$-schéma affine. L'hypothèse est satisfaite si $S$ est recouvert par
des ouverts au-dessus desquels toute partie finie de $X$ est contenue dans un
ouvert affine, en particulier si $X$ est quasi-projectif sur $S$
(corollaire 1). Le corollaire 2 ajoute que, sous ces conditions, la restriction
de Weil d'un sous-schéma fermé est un sous-schéma fermé. C'est la condition
d'existence qu'on trouve aujourd'hui dans le livre de Bosch, Lütkebohmert et
Raynaud (§ 7.6)\footnote{La page 42 appelle ces conditions « hypothèses
abracadabrantes », lecture incertaine. L'alternative (i) de la page 41 (« $T
\to S$ est \ldots{} injectif ») est barrée, et l'ordre (i)--(ii) suit les
traits de renvoi.}.

\subsection*{Ce que la restriction de Weil préserve (pages 43 à 48)}

Le foncteur $\mathrm{Res}_{T/S}$ est exact à gauche :
$\mathrm{Res}_{T/S}(T) = S$ et $\mathrm{Res}_{T/S}(X \times_Z Y) =
\mathrm{Res}_{T/S}(X) \times_{\mathrm{Res}_{T/S}(Z)} \mathrm{Res}_{T/S}(Y)$. Il
transforme immersions ouvertes en immersions ouvertes, recouvrements ouverts
en recouvrements ouverts (sous la condition du théorème), immersions fermées
en immersions fermées, morphismes affines en morphismes affines,
morphismes quasi-projectifs en morphismes quasi-projectifs, et morphismes
projectifs, ou propres, en morphismes projectifs, ou propres, \emph{pourvu que
$T \to S$ soit étale}\footnote{La page dit « pourvu que $f$ séparable », et
une marge (*) commente : il faut pouvoir descendre la propreté le long de $f$.
La restriction est nécessaire : le long d'une extension purement inséparable
non triviale, la restriction de Weil d'une variété propre de dimension $> 0$
n'est pas propre (Conrad--Gabber--Prasad, \emph{Pseudo-reductive groups},
appendice A.5, cité de mémoire). Une autre marge pose la question : si $T/S$
est fini et fidèlement plat et $Y \times_S T$ propre sur $T$, $Y$ est-il propre
sur $S$ ? « Oui dans le cas $S$ loc.\ noeth., car on utilise encore le critère
de Chow. »}.

\textbf{Proposition} (page 44). \emph{Soit $g : X \to Y$ un morphisme de
$T$-schémas et $\mathcal{L}$ un fibré en droites sur $X$ ample relativement à
$g$. Soit $\mathrm{ev} : \mathrm{Res}_{T/S}(X) \times_S T \to X$ le morphisme
universel et $p : \mathrm{Res}_{T/S}(X) \times_S T \to \mathrm{Res}_{T/S}(X)$ la
projection, finie localement libre. Alors $\det p_{*}\mathrm{ev}^{*}\mathcal{L}$,
la norme de $\mathcal{L}$, est ample relativement à $\mathrm{Res}_{T/S}(g)$.}

Il le note $\det f_{*}(\mathcal{L})$ et ajoute « (resp. très ample) »\footnote{La
version « très ample » est celle de la page, et nous ne la vérifions pas ;
la version « ample » est celle qu'on sait établir par la norme (EGA II 6.6).}.
La démonstration se ramène, la question étant locale et compatible aux
produits, au cas $X = \mathbf{P}^m_T$, $Y = T$, $S = \mathrm{Spec}\,A$, $T =
\mathrm{Spec}\,B$ avec $B$ libre de rang $\nu$ sur $A$. Il faut alors (A)
déterminer une grassmannienne et un fibré ample dessus, et (B) représenter,
sur un schéma $G$ portant une $\mathcal{O}_G$-algèbre $\mathcal{A}$ localement
libre de rang $\nu$, un $\mathcal{A}$-module libre $\mathcal{F}$ et un
sous-$\mathcal{O}_G$-module $\mathcal{F}_1$ à quotient localement libre de rang
$k\nu$, le lieu où $\mathcal{F}_1$ est un sous-$\mathcal{A}$-module et
$\mathcal{F}/\mathcal{F}_1$ un $\mathcal{A}$-module localement libre de rang $k$.
En coordonnées (pages 45 à 47) : la stabilité de $\mathcal{F}_1$ par
$\mathcal{A}$ est l'annulation de coefficients $\varphi_{i\lambda\mu}$, d'où un
sous-schéma fermé $G_1$ ; sur $G_1$, la liberté locale de $M =
\mathcal{F}/\mathcal{F}_1$ est l'inversibilité d'un des déterminants
$\det(u_\xi)$ pour $\xi \in M^k$, d'où un ouvert. Le lieu cherché est donc
localement fermé, et la condition est ponctuelle, compatible aux extensions du
corps résiduel.

\emph{Finitude} (pages 47 et 48). Si $g : X \to Y$ est de type fini et $S$
localement noethérien, $\mathrm{Res}_{T/S}(g)$ est de type fini (sous la
condition d'existence). On se ramène à $Y = T$ et l'on recouvre $X$ par un
nombre fini d'ouverts affines $U_i$ tels que les $U_i^{\nu} = U_i \times_S \cdots
\times_S U_i$ recouvrent ceux des $U^{\nu}$ (c'est possible parce que $X^{\nu}$
est un espace noethérien), puis au cas affine, et enfin, par la proposition de
la page 37, au cas d'un fibré vectoriel\footnote{La page écrit, pour ce
dernier cas, « $X = \mathrm{Spec}(\underline{F})$ \ldots{} $f_{*}(X/T) =
\mathrm{Spec}(f_{*}(\underline{F}))$ » ; on lit la formule de la page 37, seule
correcte à cette place.}.

\textbf{Remarque} (page 48). \emph{La dimension des fibres de
$\mathrm{Res}_{T/S}(X)$ n'est pas bornée par $n$ fois celle des fibres de $X/T$.}
Exemple : $S = \mathrm{Spec}\,k$, $T$ radiciel sur $S$, $X = T \times_S T$ (de
dimension relative nulle) ; alors $\mathrm{Res}_{T/S}(X) = \underline{\mathrm{Hom}}_k(T,
T)$, qui peut être de dimension $> 0$. Ainsi pour $T = \mathrm{Spec}\,k[x]/(x^2)$
en caractéristique $\neq 2$, les endomorphismes $x \mapsto a + bx$ avec $a^2 =
0$ forment, sur les points, une droite.

\subsection*{Lissité, étalitude, dimension (pages 49 à 53)}

\textbf{Proposition} (page 49). \emph{Soit $g : X \to Y$ un morphisme de type
fini de $T$-schémas dont les restrictions de Weil existent, $S$ et $T$
localement noethériens. (i) Si $g$ est lisse, $\mathrm{Res}_{T/S}(g)$ l'est.
(ii) Si $g$ est étale, $\mathrm{Res}_{T/S}(g)$ l'est.}

(i) est le critère infinitésimal : un problème de relèvement le long d'une
immersion fermée nilpotente $S'_1 \to S'$ devient, après $\times_S T$, un
problème de relèvement le long de $T'_1 \to T'$, que la lissité de $g$ résout
localement, $T'$ étant affine sur $S'$ affine. Pour (ii), le même raisonnement
donne la non-ramification ; reste la platitude, qu'on obtient en remarquant
qu'un relèvement existe point par point, $T'$ étant, au-dessus d'un corps
algébriquement clos, un nombre fini de schémas nilpotents
connexes\footnote{Lecture mince : la démonstration de (ii), page 50, est
incertaine mot à mot, et on n'en retient que l'argument d'ensemble.}.

\textbf{Dimension.} Si $g$ est lisse de dimension relative $d$ et $T/S$ de
degré $n$ constant, $\mathrm{Res}_{T/S}(g)$ est lisse de dimension relative $nd$.
Plus finement (corollaire 2, page 51), pour $s \in S$ notons $D(s)$
l'ensemble des dimensions des composantes des fibres de $X_s \to Y_s$, et
$D'(s)$ l'analogue pour $\mathrm{Res}_{T/S}(g)_s$ ; si $T_s$ est local, $D'(s) =
n D(s)$\footnote{La page écrit $D(s) = n\,D'(s)$, puis conclut « dimension $d$
pour $g$, dimension $nd$ pour $f_{*}(g)$ », ce qui demande $D'(s) = n\,D(s)$.
La démonstration (page 52) se ramène d'ailleurs au cas où $T_s$ est local :
$\mathrm{Res}$ d'une somme est le produit des $\mathrm{Res}$. Dans le cas général
l'égalité d'ensembles est fausse : si $T_s$ a deux points rationnels et que
les fibres de $g$ ont des composantes de dimension $1$ et $2$, celles de
$\mathrm{Res}(g)_s$, produits de fibres, ont aussi la dimension $3$.}. Une
première version, encadrée puis barrée (page 51), bornait la dimension par le
produit de $n$ et d'un « maximum des rangs des anneaux locaux artiniens des
$T_s$ » ; c'est le degré $[T_s : k(s)]$ qui convient\footnote{Corollaire 1 de
la page 50, lecture incertaine (« rgs », « anneaux locaux artiniens »).}.

\textbf{Corollaire 1 bis} (page 53). \emph{Si $g : X \to Y$ est fini étale de
degré $\lambda$, $\mathrm{Res}_{T/S}(g)_s$ est fini étale de degré
$\lambda^{\nu_s}$, où $\nu_s$ est le nombre de points géométriques de $T_s$ (son
degré séparable).} Une section au-dessus de $T_{\bar s}$ se choisit en effet
indépendamment en chacun des $\nu_s$ points, et se prolonge de façon unique aux
épaississements nilpotents.

\textbf{Corollaire 3} (page 53). \emph{Conditions équivalentes : (i) $\nu(s)$ est
localement constant ; (ii) $\mathrm{Res}_{T/S}$ transforme revêtements finis
étales en revêtements finis étales ; (iii) $\mathrm{Res}_{T/S}(T')$ est fini étale
pour tout $T'$ fini étale sur $T$ ; (iv) $\mathrm{Res}_{T/S}(T \sqcup T)$ est fini
étale sur $S$.} Une marge relie (i) $\Rightarrow$ (ii) $\Rightarrow$ (iii)
$\Rightarrow$ (iv) $\Rightarrow$ (i).

\textbf{Corollaire 4} (page 54). \emph{Si $g$ est lisse et surjectif,
$\mathrm{Res}_{T/S}(g)$ est surjectif.}

\subsection*{Torseurs, et les composantes connexes relatives (pages 54 à 58)}

La restriction de Weil transforme torseurs en torseurs : un revêtement
principal de groupe fini $G$ de $Y$ devient un torseur sur
$\mathrm{Res}_{T/S}(Y)$ sous le groupe $\mathrm{Res}_{T/S}(G_T)$\footnote{Pages 54 et
55, lecture mince : la proposition générale est en grande partie incertaine,
et une marge conseille de prendre « plutôt la notion de fibré principal
homogène ».}. D'où la question de déterminer, pour un ensemble fini $I$, le
schéma étale $\mathrm{Res}_{T/S}(T \times I)$.

La réponse occupe les pages 55 à 58. Supposons $\nu(s)$ localement constant,
comme au corollaire 3\footnote{La page ne le suppose pas à cet endroit ; sans
cette condition $\mathrm{Res}_{T/S}(T \times I)$ est étale mais pas fini, et le
revêtement $A(T)$ ci-dessous n'existe pas. La page dit aussi supposer $S$
connexe, « ce qu'on devrait pas être nécessaire ».}. Le foncteur $I \mapsto
\mathrm{Res}_{T/S}(T \times I)$, des revêtements finis étales de $S$ dans
eux-mêmes, est exact à gauche ; au-dessus d'un point géométrique $a$ de $S$
connexe, c'est un foncteur exact à gauche sur les $\pi_1(S,a)$-ensembles finis,
donc de la forme $E \mapsto \mathrm{Hom}_{\pi_1}(A, E)$ pour un $\pi_1$-ensemble
fini $A = A(T)$. Le revêtement étale correspondant, qu'il note
$f^{\mathrm{\acute{e}t}}_{*}(T)$, est caractérisé par
\[
  \mathrm{Hom}_S\bigl(f^{\mathrm{\acute{e}t}}_{*}(T), S'\bigr) \simeq \mathrm{Hom}_S(T, S')
\]
pour tout revêtement fini étale $S'$ de $S$ ; d'où, pour $S' =
f^{\mathrm{\acute{e}t}}_{*}(T)$, un morphisme canonique $T \to
f^{\mathrm{\acute{e}t}}_{*}(T)$, surjectif et à fibres géométriques connexes. Une
réciproque : un morphisme surjectif à fibres géométriques connexes de $T$ sur
un revêtement fini étale $\overline{T}$ identifie $\overline{T}$ à
$f^{\mathrm{\acute{e}t}}_{*}(T)$\footnote{La page 58 énonce la réciproque pour un
morphisme « surjectif et fini injectif », lecture incertaine, et sa dernière
ligne n'est pas déchiffrée. On donne la réciproque juste.}.

C'est le schéma des composantes connexes relatives $\pi_0(T/S)$, et la
variance demande un mot. Le nom est celui d'une image directe, mais
$f^{\mathrm{\acute{e}t}}_{*}$ est un adjoint \emph{à gauche} (de l'inclusion des
revêtements finis étales dans les $S$-schémas), alors que la restriction de
Weil est un adjoint à droite. Les deux se rejoignent par la formule même qui
le définit : $\mathrm{Res}_{T/S}(T \times I)$ est le schéma des applications de
$\pi_0(T/S)$ dans $I$\footnote{Le nom $\pi_0(T/S)$ et la référence (EGA IV 6.8
pour les cas finis ; D.~Romagny, « Composantes connexes et irréductibles en
familles », 2011, cité de mémoire) sont les nôtres.}.

\pagerange{60}{79}
\section*{VII. Un foncteur de nature locale et son adjoint : Weil et Greenberg (pages 60 à 79)}

\subsection*{Le cadre, et les exemples (pages 60 à 66)}

La page 60 ouvre une nouvelle rédaction, « (1) Set up optimal »\footnote{Le
titre est souligné et sa lecture très incertaine.}. Soit $\mathbb{W}$ un
foncteur covariant des $S$-schémas dans des schémas, \emph{de nature locale} :
il transforme immersions ouvertes en immersions ouvertes, et commute aux
réunions et aux intersections d'ouverts. Exemples :

\begin{itemize}
  \item[I.] Le changement de base $\mathbb{W}_{S'}(T) = T \times_S S'$, pour un
  $S$-schéma $S'$ ; I$'$ : $S'$ fini localement libre sur $S$.
  \item[II.] Un schéma en anneaux commutatifs $\Lambda$ sur $S$ tel que, pour
  tout $T$, le faisceau d'anneaux $\Gamma(\Lambda_T/T)$ sur l'espace $|T|$ fasse de
  $|T|$ un schéma $\mathbb{W}(T)$ : $|\mathbb{W}(T)| = |T|$, et les anneaux
  locaux sont « tordus » par $\Lambda$. C'est le cas des vecteurs de Witt
  tronqués, et ce sont les foncteurs de Greenberg ; II$'$ : $\mathbb{W}(T) \to
  T$ surjectif à noyau nilpotent sur les faisceaux d'anneaux.
  \item[III.] $S' = \mathrm{Spec}\,\mathcal{A}$, $\mathcal{A}$ localement libre,
  $S' \to S$ homéomorphisme universel, et $\Lambda = \mathbf{V}(\check{\mathcal{A}})$,
  à la fois I et II.
\end{itemize}

Pour $\mathfrak{X}$ un schéma sur $\mathbb{W}(S)$, on considère le foncteur sur
les $S$-schémas $T \mapsto \mathrm{Hom}_{\mathbb{W}(S)}(\mathbb{W}(T),
\mathfrak{X})$, et l'on note $\mathbb{W}^{a}(\mathfrak{X})$ le schéma qui le
représente quand il existe :
\[
  \mathrm{Hom}_S\bigl(T, \mathbb{W}^{a}(\mathfrak{X})\bigr) \simeq
  \mathrm{Hom}_{\mathbb{W}(S)}\bigl(\mathbb{W}(T), \mathfrak{X}\bigr).
\]
$\mathbb{W}^{a}$ est donc un adjoint \emph{à droite} de $\mathbb{W}$, défini sur
la sous-catégorie des $\mathfrak{X}$ pour lesquels il existe\footnote{La page
dit « foncteur adjoint (au sens de Kan) » sans préciser le côté ; la formule
le fixe. Celle qui clôt la page 62, $\mathrm{Hom}_{S^{a}}(T^{a}, \mathfrak{X})
\simeq \mathrm{Hom}_S(T, \mathfrak{X})$, est coupée au bord et ses exposants
sont incertains ; on ne la retient pas.}. Dans l'exemple I c'est la
restriction de Weil $\mathrm{Res}_{S'/S}$, dans l'exemple II le foncteur de
Greenberg\footnote{M.~J.~Greenberg, « Schemata over local rings », \emph{Ann.
of Math.} 73 (1961), cité de mémoire. Le dossier 23 (pages 19 à 22) rencontre
la même construction, le long d'un épaississement infinitésimal, sous la
notation $X^{\{\Lambda\}}$ ; le dossier 43 (pages 111 à 113) l'emploie pour
présenter les schémas en groupes finis comme noyaux d'isogénies.}\footnote{Un
cadre axiomatique commun aux foncteurs de Greenberg et aux restrictions de
Weil (logarithmiques) figure dans I.~Vanni, \emph{Greenberg functors and jumps
of tori}, thèse, Sapienza Università di Roma, 2025, au moyen de « catégories
géométriques », avec représentabilité et conservation de la quasi-compacité et
de la quasi-séparation ; nous n'en avons lu que le résumé, et ne savons donc
pas si ce cadre est celui des pages 60 à 79. Que le foncteur de
Greenberg coïncide avec la restriction de Weil en égale caractéristique est
dans A.~Bertapelle et C.~D.~González-Avilés, « The Greenberg functor
revisited », \emph{Eur. J. Math.} 4 (2018), remarque 7.17 (c) et remarque
11.2 (arXiv:1311.0051, v6) : c'est, dans un cas particulier, le rapprochement
des exemples I et II.}.

Propriétés (pages 63 à 66) : $\mathbb{W}^{a}$ commute aux limites projectives
finies, donc transforme schémas en groupes, en anneaux, en schémas du même
type ; il commute au changement de base, $\mathbb{W}^{a}(\mathfrak{X}) \times_S T
= \mathbb{W}^{a}_T(\mathfrak{X} \times_{\mathbb{W}(S)} \mathbb{W}(T))$. Si
$\mathbb{W}^{a}(\mathbf{A}^1_{\mathbb{W}(S)})$ existe, c'est un schéma en anneaux
sur $S$ qui \emph{engendre} $\mathbb{W}$ : c'est $\Lambda$ dans l'exemple II, et
$\mathbf{V}(\check{\mathcal{A}})$ dans l'exemple I$'$. Alors, pour $\mathfrak{X}$
fermé dans $\mathbf{A}^r_{\mathbb{W}(S)}$ défini par $F_1, \ldots, F_s$, les $F_j$
ont des coefficients dans $\Gamma(\Lambda/S)$, définissent $\Lambda^r \to
\Lambda^s$, et $\mathbb{W}^{a}(\mathfrak{X})$ est l'image inverse de la section
nulle. Exemples donnés (page 65) : $\mathbb{W}^{a}$ de $\mathbf{G}_a^n$, de
$\mathrm{GL}_n$ sur $\mathbb{W}(S)$. Enfin un morphisme $\mathbb{W}' \to \mathbb{W}$ de tels foncteurs induit
$\mathbb{W}^{a}(\mathfrak{X}) \to \mathbb{W}'^{a}(\mathfrak{X} \times_{\mathbb{W}(S)}
\mathbb{W}'(S))$, et l'expression de droite est un bifoncteur covariant en
$\mathfrak{X}$, contravariant en $\mathbb{W}'$ ; la variance est bien celle de
la page.

\subsection*{Le cas affine (pages 66 à 68)}

Supposons que $\mathbb{W}$ transforme schémas affines en schémas affines ; sur
un ouvert affine $\mathrm{Spec}\,A$ de $S$, il est alors donné par $\mathbb{W}(B)
= \mathrm{Hom}_{A\text{-alg}}(\Lambda, B)$, $\Lambda$ l'anneau d'un schéma en
anneaux affine sur $S$ (cas I$'$, et cas II quand $\Lambda$ est affine).

\textbf{Proposition.} \emph{Si $\mathfrak{X}$ est affine sur $\mathbb{W}(S)$,
$\mathbb{W}^{a}(\mathfrak{X})$ existe et est affine sur $S$. Si l'anneau de
$\mathfrak{X}$ est engendré par $r$ éléments et celui de $\Lambda$ par $s$
éléments, celui de $\mathbb{W}^{a}(\mathfrak{X})$ l'est par $rs$
éléments.}\footnote{L'énoncé du corollaire (pages 67 et 68) est en partie
illisible ; on reconstruit la seule lecture compatible avec la démonstration,
qui plonge $\mathbb{W}^{a}(\mathfrak{X})$ dans $\Lambda^r$.}

Pour $\mathfrak{X} = \mathrm{Spec}\,\mathcal{C}$ quelconque, on écrit $\mathcal{C} =
\varinjlim \mathcal{C}_i$ (sous-algèbres de type fini) et l'on trouve
$\mathbb{W}^{a}(\mathfrak{X}) = \mathrm{Spec}\,\mathcal{D}$, avec $\mathcal{D} =
\varinjlim \mathcal{D}_i$, $\mathcal{D}_i$ l'anneau de
$\mathbb{W}^{a}(\mathrm{Spec}\,\mathcal{C}_i)$.

\subsection*{Recollement (pages 69 à 76)}

Sans hypothèse sur $\mathbb{W}$, deux conditions sur un recouvrement ouvert
$\mathfrak{X} = \bigcup \mathfrak{X}_i$ :
\begin{itemize}
  \item[($\alpha$)] (lemme 1) si $\mathbb{W}^{a}(\mathfrak{X}) = X$ existe et
  $\mathfrak{U} \subset \mathfrak{X}$ est ouvert, d'image inverse $\mathfrak{U}'$
  dans $\mathbb{W}(X)$, il existe un ouvert $U$ de $X$ tel qu'un morphisme
  $\mathbb{W}(T) \to \mathbb{W}(X)$ se factorise par $\mathfrak{U}'$ si et
  seulement s'il se factorise par $\mathbb{W}(U)$ ; alors $U =
  \mathbb{W}^{a}(\mathfrak{U})$\footnote{La page écrit « est un
  $\mathbb{W}(\mathfrak{U}')$ » ; on attend $\mathbb{W}^{a}(\mathfrak{U})$, et c'est
  ce que la suite emploie.} ;
  \item[($\beta$)] (lemme 2) pour tout $\mathbb{W}(T) \to \mathfrak{X}$, il existe
  un recouvrement ouvert $T = \bigcup T_i$ avec $\mathbb{W}(T_i)$ envoyé dans
  $\mathfrak{X}_i$.
\end{itemize}
Sous ($\beta$), et si les $\mathbb{W}^{a}(\mathfrak{X}_i)$ existent avec leurs
intersections ouvertes, $\mathbb{W}^{a}(\mathfrak{X})$ existe et est leur
recollement ; la démonstration (pages 70 à 73) vérifie la condition de cocycle
et l'universalité, par raffinement de deux recouvrements\footnote{Page 72, la
factorisation finit sur la page en « $X_i$ » ; on lit $\mathfrak{X}_i$.}. La
condition ($\beta$) est automatique si $|\mathbb{W}(T)| \simeq |T|$ (cas II$'$
et III, lemme 3) ; ($\alpha$) l'est dans les cas I$'$ et « Greenberg
spécial » (lemme 4 ; pour I$'$, parce qu'un morphisme fini est fermé). Un corollaire
(page 74) note que les conditions sont encore satisfaites dans le cas I quand
$S'$ est propre sur $S$ et $\mathbb{W}^{a}(\mathfrak{X})$ localement noethérien ;
une ligne suivante, « Mauvaises \ldots{} I$'$ », est presque illisible.

\textbf{Théorème 1} (page 74). \emph{Pour un foncteur de Greenberg spécial, ou
pour le changement de base le long de $S' \to S$ fini, localement libre et
radiciel, $\mathbb{W}^{a}(\mathfrak{X})$ existe pour tout $\mathfrak{X}$ sur
$\mathbb{W}(S)$, et $\mathbb{W}^{a}$ est de nature locale.}

C'est exact : dans ces deux cas $|\mathbb{W}(T)| = |T|$, et une section au-dessus
d'un point a son image en un seul point, donc dans un ouvert
affine\footnote{« Greenberg spécial » n'est pas défini sur les pages ; c'est
vraisemblablement le cas II avec $\Lambda$ affine, celui du lemme 4.}\footnote{Les
deux cas du théorème sont publiés. Pour le foncteur de Greenberg :
A.~Bertapelle et C.~D.~González-Avilés, « The Greenberg functor revisited »,
\emph{Eur. J. Math.} 4 (2018), § 7 : le foncteur $h_R$ est local pour la
topologie de Zariski et respecte immersions ouvertes et fermées (proposition
7.4 et (7.8)), et la proposition-définition 7.10 montre que $\mathrm{Gr}^R(Z)$
existe pour tout $R$-schéma $Z$, avec $\mathrm{Hom}_k(Y, \mathrm{Gr}^R(Z)) \simeq
\mathrm{Hom}_R(h_R(Y), Z)$ (7.12) : c'est le même énoncé, $\mathrm{Gr}^R$ adjoint à
droite de $h_R$ ; le § 17 (vers (17.11)) l'étend aux schémas formels adiques.
Pour la restriction de Weil le long d'un $S' \to S$ fini localement libre qui
est un homéomorphisme universel, l'existence pour tout $X$ est le corollaire
A.5 du même article, repris dans O.~Overkamp, « Néron models of
pseudo-Abelian varieties », arXiv:2003.03764 (v2, 2021), proposition 2.1, qui
s'appuie pour le cas affine et le recollement le long d'immersions ouvertes
sur Bosch--Lütkebohmert--Raynaud, § 7.6, proposition 2 (i) et théorème 4 (tel
que cité là) : même énoncé que le second cas de la page 74.}.

\textbf{Théorème 1 bis} (pages 74 et 75). \emph{Dans le cas I$'$ : (i) si
$\mathrm{Res}_{S'/S}(\mathfrak{X})$ existe, $\mathrm{Res}_{S'/S}(\mathfrak{U})$ existe
pour tout ouvert $\mathfrak{U}$ et en est un ouvert ; (ii) si $\mathfrak{X} =
\bigcup \mathfrak{X}_i$ et que, pour tout $s$, toute section de $\mathfrak{X}_s$
sur $S'_s$ définie sur une extension de $k(s)$ tombe dans un $\mathfrak{X}_i$,
l'existence des $\mathrm{Res}_{S'/S}(\mathfrak{X}_i)$ entraîne celle de
$\mathrm{Res}_{S'/S}(\mathfrak{X})$, réunion des
précédents.}\footnote{L'énoncé (ii) porte deux mots illisibles à la place de
« section » ; on reprend la formulation de la proposition 2 bis de la page
39, que la démonstration de la page 75 suit.} Ce sont, dans le nouveau
cadre, la proposition 2 bis et le théorème des pages 39 à 41 ; les
corollaires (pages 75 et 76) redonnent l'existence quand toute partie finie
des fibres est dans un ouvert affine, et la finitude par le même argument
noethérien sur $\mathfrak{X} \times_S \cdots \times_S \mathfrak{X}$.

\subsection*{Immersions, morphismes de type fini, affines (pages 77 à 79)}

Pour $\mathbb{W}$ de type I$'$ ou Greenberg spécial : $\mathbb{W}^{a}$ transforme
immersions fermées en immersions fermées (dans le cas I$'$, sous la condition
d'existence pour le but), immersions en immersions si le but est
quasi-projectif, morphismes séparés en morphismes séparés (appliquer ce qui
précède à la diagonale), morphismes de type fini en morphismes de type fini,
et morphismes affines en morphismes affines\footnote{La dernière proposition,
page 79, n'a que son énoncé : « Comme dessus, pour un morphisme
\emph{affine} ». Page 77, le « $S^{*}$ » de « quasi-projectif sur $S^{*}$ » est
sur la page, sans explication.}\footnote{Pour le foncteur de Greenberg, ces
propriétés de permanence sont publiées : A.~Bertapelle et
C.~D.~González-Avilés, « The Greenberg functor revisited », \emph{Eur. J.
Math.} 4 (2018), proposition 11.3 ($\mathrm{Gr}^R$ conserve les morphismes
quasi-compacts, quasi-séparés, séparés, localement de type fini, de type fini
et affines) et remarque 7.17 (b) (immersions fermées) ; c'est la même liste,
établie pour ce seul cas, celui de la restriction de Weil étant traité à part
(voir la note au théorème 1).}.

\pagerange{80}{121}
\section*{VIII. « Schémas localement compacts » (pages 80 à 121)}

La couverture de la page 80 porte ce titre, sous un premier titre biffé
(« Notions d'espaces \ldots{} »). Le mot « compact » y a le sens de Bourbaki,
séparé compris, et il s'agit de regarder un schéma à travers ses
parties constructibles seulement, c'est-à-dire en oubliant tout ce que les
morphismes ont de « continu » pour ne garder que ce qu'ils ont de « par
morceaux ».

\subsection*{Applications constructibles et ensembles schématiques (pages 81 à 100)}

Les pages paires portent les notes, les impaires sont des versos sans
rapport, et la main est rapide : on en donne les définitions et les énoncés,
la prose qui les relie étant souvent incertaine.

\emph{Constructibilité locale} (pages 81 et 83). Une partie $Z$ est constructible
si et seulement si elle est quasi-compacte et localement constructible ; une
partie localement fermée $Z$ d'un schéma quasi-compact et quasi-séparé est
constructible si et seulement si elle est quasi-compacte et localement définie,
dans un ouvert affine, par un nombre fini d'équations\footnote{Page 83 commence
une liste d'équivalences, pour $T$ fermé irréductible, entre l'existence d'une
partie constructible dense dans $T$ et une condition (ii) que le bas de la
feuille, déchiré, interrompt.}.

\emph{Schémas atomiques} (page 86). Un schéma \emph{atomique} est une somme de
spectres de corps ; ils forment une sous-catégorie pleine, qui se plonge
pleinement fidèlement dans les foncteurs covariants des corps vers les
ensembles (par $\coprod_x \mathrm{Spec}\,k(x) \mapsto \coprod_x \mathrm{Hom}(k(x),
-)$, et le lemme de Yoneda). Tout schéma $X$ définit le schéma atomique $|X| =
\coprod_{x \in X} \mathrm{Spec}\,k(x)$, et $X \mapsto |X|$ est fidèle sur les
schémas réduits.

\textbf{Définition.} \emph{Une application constructible $f : X \to Y$ est un
morphisme $|X| \to |Y|$ tel que tout ouvert affine $U$ de $X$ admette une
partition finie en parties localement fermées constructibles $U_i$, sur
chacune desquelles $f$ provient d'un morphisme $(U_i)_{\mathrm{r\acute ed}} \to Y$.
Un ensemble schématique est un schéma quasi-compact et quasi-séparé, vu comme
objet de la catégorie dont les morphismes sont les applications
constructibles.}\footnote{Le mot est lu « schématique » sur les tuiles à
700 dpi ; le premier passage lisait « élémentaire ». Il revient jusqu'à la
page 100. Le nom n'a pas d'équivalent courant : ces objets sont ce que la
théorie des modèles appelle des ensembles définissables, et c'est aussi le
cadre de l'anneau de Grothendieck des variétés, où deux variétés découpables
en morceaux isomorphes sont identifiées. Le rapprochement est le nôtre.}

Il les étudie (pages 88 à 94). Un lemme, cerné en marge : sur un schéma
quasi-compact quasi-séparé, tout recouvrement fini par des parties
constructibles est raffiné par une partition finie en parties localement
fermées constructibles affines. Puis : pour $X$ quasi-compact quasi-séparé, $f$ est
constructible si et seulement s'il existe une partition finie de $X$ en
parties localement fermées constructibles affines sur lesquelles $f$ est
morphique ; l'image inverse d'une partie constructible par une application
constructible est constructible, de sorte que la notion de partie
constructible est invariante par isomorphisme d'ensembles schématiques ; $f$ est un isomorphisme d'ensembles
schématiques si et seulement s'il existe des partitions finies $(X_i)$,
$(Y_i)$ en parties localement fermées constructibles avec $f$ induisant des
isomorphismes $X_{i,\mathrm{r\acute ed}} \simeq Y_{i,\mathrm{r\acute ed}}$. D'où :
\begin{itemize}
  \item tout ensemble schématique est isomorphe à un $|X|$ avec $X$ affine
  (somme finie des morceaux) ;
  \item si $X$ est noethérien, l'espace sous-jacent de $Y$ l'est\footnote{La
  page dit « Cor.~0. Si $X$ noethérien, $Y$ l'est ». C'est vrai de l'espace
  sous-jacent, réunion finie de morceaux noethériens, non du schéma $Y$, qui
  peut porter des nilpotents non noethériens que les applications
  constructibles ne voient pas.} ;
  \item tout morphisme d'ensembles schématiques est isomorphe à un morphisme de
  schémas, vu comme application constructible\footnote{Corollaire 2 de la
  page 92. Un corollaire 3, en marge, est perdu dans des trous du papier ;
  la même marge parle d'« ens.\ schématique noethérien ».} ;
  \item les morphismes de type fini, les produits fibrés et les limites
  projectives finies existent dans la catégorie, et le foncteur des schémas
  quasi-compacts quasi-séparés vers elle commute aux produits fibrés.
\end{itemize}

Les pages 96 et 98 posent une question : pour un système projectif à
transitions affines, a-t-on $\mathrm{Cons}(\varprojlim X_i) \simeq \varprojlim
\mathrm{Cons}(X_i)$ ? Les formules écrites interchangent une limite inductive
sur les partitions finies et une limite projective sur $i$, et la page ne
conclut pas (« ??? »). Elle avance ensuite que les ensembles
« quasi-schématiques » de type fini sur un noethérien forment une « catégorie
géométrique (?) », et la page 100 en dresse la liste d'axiomes : limites
projectives finies, quotients par les groupoïdes d'équivalence, épimorphismes
effectifs et universels, et leurs duaux ; puis l'« axiome géométrique » $X
\times \coprod_i Y_i \simeq \coprod_i X \times Y_i$. La liste rappelle les axiomes
de Giraud des topos (SGA 4, 1963--1964), transposés à une catégorie
d'espaces ; la page s'arrête sur « $R \to S$ »\footnote{Le rapprochement avec
les axiomes de Giraud est le nôtre. En tête de la liste : « Tout épimorphisme
est un morphisme de descente », avec « application surjective » biffé et
remplacé par « morphisme fid.\ plat ».}.

\subsection*{La topologie constructible (pages 103 à 110)}

\textbf{Définition 1.} \emph{La topologie constructible d'un espace $X$ est celle
dont les ouverts sont les réunions de parties localement constructibles ; on
note $X^{\mathrm{cons}}$ l'espace ainsi topologisé. Une partie pro-constructible
est un fermé de $X^{\mathrm{cons}}$, c'est-à-dire une intersection de parties
localement constructibles.}\footnote{La page dit « partie analytique ». Le
nom d'aujourd'hui, « pro-constructible », est celui d'EGA IV 1.9 (1964), où
la topologie constructible est aussi introduite (1.9.13). On ne sait pas
d'après le dossier lequel précède l'autre.}

\textbf{Proposition 1.} \emph{Si $X$ a une base d'ouverts rétrocompacts :
(0) les parties constructibles forment une base de $X^{\mathrm{cons}}$ ; (i)
$X^{\mathrm{cons}} \to X$ est continue ; (ii) $X^{\mathrm{cons}}$ est séparé (et
totalement discontinu) si et seulement si $X$ est $T_0$ ; (iii) sur une
partie constructible $E$, la topologie constructible de $E$ est induite par
celle de $X$.}

Un lemme (pages 104 et 105) établit la transitivité : pour $E$ localement
constructible, une partie de $E$ est (localement) constructible dans $E$ si et
seulement si elle l'est dans $X$ (EGA $0_{\mathrm{III}}$ 9.1.8 et 9.1.13,
auxquels la page renvoie).

\textbf{Corollaires} (pages 105 à 107). \emph{Si de plus les ouverts de base
sont quasi-compacts pour la topologie constructible et $X$ est $T_0$ :
$X^{\mathrm{cons}}$ est localement compact, et compact si $X$ est quasi-compact ;
une partie est fermée dans $X^{\mathrm{cons}}$ si et seulement si elle y est
compacte ; $E$ est constructible si et seulement si $E$ et $X - E$ sont
pro-constructibles, si et seulement si $E$ est ouvert et fermé dans
$X^{\mathrm{cons}}$.}

Une parenthèse de la page 108 rappelle le critère topologique dont tout
cela dépend : si $X$ a une base d'ouverts quasi-compacts $U_i$, une application
continue $f : X' \to X$ est quasi-compacte si et seulement si les
$f^{-1}(U_i)$ le sont ; et les $U_i$ sont rétrocompacts si et seulement si tout
ouvert quasi-compact l'est.

\textbf{Proposition 2 et corollaires} (pages 107 à 109). \emph{Si $f : X \to Y$
est continue et que l'image inverse d'un ouvert rétrocompact soit
rétrocompacte, $f$ ramène les parties (localement) constructibles en parties
(localement) constructibles, et $f^{\mathrm{cons}}$ est continue. Pour $X$, $Y$ de
la classe précédente : $f$ est quasi-compacte si et seulement si
$f^{\mathrm{cons}}$ est propre ; si $X$ et $Y$ sont quasi-compacts, $f$ envoie
constructibles sur constructibles si et seulement si $f^{\mathrm{cons}}$ est
ouverte, ce qui a lieu si $f$ est bijective ; $f(X)$ est pro-constructible ;
une partie $E$ de $f(X)$ est pro-constructible si et seulement si $f^{-1}(E)$
l'est, et si $f$ est surjective, $E$ est constructible si et seulement si
$f^{-1}(E)$ l'est.}\footnote{Page 109, (iii) écrit « $f^{-1}(Y)$ » ; on lit
$f^{-1}(E)$. Une marge renvoie à Bourbaki, \emph{Topologie générale}, I, § 10,
pour les applications propres. Une note marginale page 107 juge l'hypothèse de
la proposition 2 « trop restrictive pour un morphisme de préschémas ».}

\textbf{Théorème} (page 110). \emph{Si $X$ est affine, $X^{\mathrm{cons}}$ est
compact. Si $X$ est quasi-séparé, $X$ est $T_0$ et a une base d'ouverts
rétrocompacts et compacts pour la topologie constructible, de sorte que
$X^{\mathrm{cons}}$ est localement compact. Si $X$ est quasi-compact et
quasi-séparé : $X^{\mathrm{cons}}$ est compact ; les parties constructibles sont
les ouverts fermés de $X^{\mathrm{cons}}$ ; les parties pro-constructibles sont
les images des morphismes quasi-compacts $X' \to X$, et l'on peut prendre
$X' \to X$ un monomorphisme.}\footnote{La page écrit « [quasi-compact et
quasi-séparé] quasi-noethérien » : un mot de plus que la parenthèse, sans
définition ; nous gardons la parenthèse. Le point (iv) de la page, qui reprend
le corollaire 3, est en partie illisible. Ces énoncés sont ceux d'EGA IV 1.9 ;
la topologie constructible fait de $X$ ce qu'on appelle aujourd'hui un espace
spectral muni de sa topologie « patch ».}

Le dossier 23 (page 72) rencontre la même topologie par l'autre bout : il
caractérise les schémas pour lesquels $X^{\mathrm{cons}} \to X$ est un
homéomorphisme (les schémas de dimension zéro) et, parmi eux, ceux dont tout
module est plat.

\subsection*{Le foncteur $\mathrm{Cons}$ et l'anneau $A^{\mathrm{cons}}$ (pages 112 à 115)}

La page 112 étend la notion d'application constructible aux couples $(X, E)$,
$E$ partie localement constructible d'un schéma $X$, et introduit la catégorie
$(\mathrm{Cons\,Aff})$ : objets les schémas affines (ou quasi-compacts
quasi-séparés), morphismes les applications constructibles. À $(X,E)$ elle
associe le préfaisceau $\mathrm{Cons}(X,E) : Y \mapsto \mathrm{App\,cons}(Y, (X,E))$
sur $(\mathrm{Cons\,Aff})$, et la page 113 affirme que $(X,E) \mapsto
\mathrm{Cons}(X,E)$ est pleinement fidèle, de même que les foncteurs
$(\mathrm{Cons\,Aff}) \to (\text{couples}) \to \mathrm{Ind}(\mathrm{Cons\,Aff})' \to
\underline{\mathrm{Hom}}((\mathrm{Cons\,Aff})^{\circ}, \mathrm{Ens})$\footnote{Le
$\mathrm{Ind}(\ldots)'$ est glosé « systèmes inductifs injectifs\ldots ». La
pleine fidélité du dernier est le lemme de Yoneda ; celle des autres est
affirmée sans démonstration.}.

\textbf{La construction de la page 114.} Soit $A$ un anneau, $X =
\mathrm{Spec}\,A$. Pour une partition finie $\mathfrak{P} = (X_i)$ de $X$ en parties
localement fermées constructibles affines, posons $A_{\mathfrak{P}} = \prod_i
\Gamma(X_{i,\mathrm{r\acute ed}}, \mathcal{O})$, et
\[
  A^{\mathrm{cons}} = \varinjlim_{\mathfrak{P}} A_{\mathfrak{P}} .
\]
Les transitions sont injectives. $\mathrm{Spec}\,A_{\mathfrak{P}} = \coprod_i
X_{i,\mathrm{r\acute ed}}$ est en bijection avec $X$, les transitions sont des
bijections, et $\mathrm{Spec}\,A^{\mathrm{cons}} = \varprojlim
\mathrm{Spec}\,A_{\mathfrak{P}}$ est donc en bijection avec $X$ ; la topologie
limite est engendrée par les ouverts des morceaux, qui sont constructibles, et
c'est la topologie constructible : $\mathrm{Spec}\,A^{\mathrm{cons}} \simeq
X^{\mathrm{cons}}$\footnote{La page dit seulement que la bijection « est
compatible avec les topologies \ldots{} comme on voit aisément sur la
description d'une base d'ouverts » ; l'argument qu'on donne est celui-là. Si
$A$ est noethérien, les anneaux locaux de $A^{\mathrm{cons}}$ sont les corps
résiduels de $A$ (on peut prendre pour morceau contenant $x$ un ouvert de
$\overline{\{x\}}$), donc $A^{\mathrm{cons}}$ est absolument plat : à
rapprocher de l'anneau absolument plat universel d'Olivier (1968) et de la
page 72 du dossier 23. Ce rapprochement est le nôtre.}\footnote{La
construction est publiée : The Stacks Project, chapitre « Pro-étale
Cohomology », section « Constructing w-local affine schemes ». Les strates
$Z(E', E'') = D(\prod E') \cap V(E'')$, $E \subset A$ fini, raffinent toute
stratification constructible finie (lemme 096Y), et la construction qui
précède le lemme 0975 donne $A_w = \varinjlim_E A_E$ avec le sous-schéma fermé
$Z = \varprojlim_E Z_E$, $Z_E = \coprod Z(E', E'')$ ; donc $Z =
\mathrm{Spec}(\varinjlim_E \prod \Gamma(Z(E', E''), \mathcal{O}))$, c'est-à-dire
$A^{\mathrm{cons}}$ au choix près d'un système cofinal de partitions. Le lemme
0975 (2)--(4) dit que $Z \to \mathrm{Spec}\,A$ est bijectif, que $Z$ est
l'ensemble des points fermés de $X_w$ et qu'il est réduit ; la définition 096A
et la remarque qui la suit, que les points fermés d'un espace w-local forment
un espace profini. L'homéomorphisme avec la topologie constructible s'en
déduit en un pas (chaque strate se relève en un ouvert fermé, et une
bijection continue d'un compact sur un séparé est un homéomorphisme), mais
n'y est pas écrit en ces termes ; et comme $Z$ est réduit d'espace profini,
$A^{\mathrm{cons}}$ est absolument plat pour tout anneau $A$, pas seulement
noethérien. La source de cette construction est B.~Bhatt et P.~Scholze,
« The pro-étale topology for schemes », \emph{Astérisque} 369 (2015), § 2.2,
que nous n'avons pas lu. Que l'algèbre absolument plate universelle d'Olivier
ait pour spectre $\mathrm{Spec}\,A$ muni de la topologie constructible, nous ne
le tenons que d'une source secondaire, P.~Belmans, « Absolute
flatification » (blog, 25 décembre 2015) ; le travail d'Olivier n'a pas été
consulté.}. C'est la même limite
sur les partitions constructibles finies que le dossier 43 (pages 91 à 97)
emploie pour définir le foncteur qu'il appelle « de Néron ».

La page 115, barrée entière, est un premier essai de la définition 1.

\subsection*{Le plan de la suite, et les schémas compacts (pages 117 à 121)}

La page 117 est un plan : N° 1 schémas compacts et localement compacts,
anneaux flasques ; N° 2 proschémas flasques, espaces schématiques compacts ;
N° 3 espaces schématiques ; N° 4 opérations, et limites projectives ; N° 5
espaces quasi-schématiques ; N° 6 « Pb d'espaces. Espaces de Néron ». Il est
suivi de deux diagrammes de catégories (pages 117 et 119) : schémas
quasi-compacts quasi-séparés, pro-schémas « flasques », schémas compacts
réduits, et leurs versions $\mathrm{Ind}$, avec des foncteurs $\gamma$, $\beta$ et
des adjoints $\gamma'$, $\beta'$ pleinement fidèles, $\mathrm{Hom}(K, \gamma X)
\simeq \mathrm{Hom}(\gamma' K, X)$, où $\gamma'$ est donc adjoint \emph{à gauche}
de $\gamma$\footnote{Plusieurs étiquettes de ces diagrammes sont incertaines
(« surjectif », « subjectif », « immersion ») et le départ de l'arc
$\overline{\gamma}$ de la page 119 n'est pas sûr. On ne rapporte que les
formules d'adjonction, qui sont lisibles.}. Ce que sont les « espaces de
Néron » n'est dit nulle part.

\textbf{Schémas compacts} (page 121). \emph{Soit $X$ un schéma compact, c'est-à-dire
quasi-compact et séparé comme espace topologique. Alors (i) $X$ est totalement
discontinu, tout point ayant un système fondamental de voisinages ouverts et
fermés ; (ii) $X$ est affine ; (iii) ses parties constructibles sont ses parties
ouvertes et fermées.} Pour (ii) : on recouvre $X$ par des ouverts affines
$X_1, \ldots, X_n$, qui sont ouverts et fermés ; les $X'_i = X_i - \bigcup_{j<i}
X_j$ sont ouverts et fermés dans des affines, donc affines, et $X$ est leur
somme. Ces schémas sont ceux de dimension zéro à espace quasi-compact, dont le
dossier 23 (page 72) donne d'autres caractérisations.

\pagerange{123}{141}
\section*{IX. « Greenberg -- Néron » (pages 123 à 141)}

La page 123 ne porte que ce titre. Ce qui suit est la partie la plus
fragmentaire du dossier : des définitions et des énoncés sans démonstration,
puis des calculs.

\subsection*{Points épais d'une limite projective (pages 125 et 127)}

Soit $(X_i)$ un système projectif de schémas essentiellement affine (les
transitions sont affines à partir d'un rang), de sorte que $X = \varprojlim X_i$
existe et que son espace est la limite des espaces. Un point $x = (x_i)$ est
\emph{épais} si
\begin{quote}
(Ep) il existe $i$ tel que pour tout $j \geq i$, $x_j$ soit un point maximal de
la fibre $\varphi_{ij}^{-1}(x_i)$.
\end{quote}
Variantes : (Ep)$'$ il existe une partie ouverte non vide $U_i$ de
$\overline{\{x_i\}}$ avec $\varphi_{ij}^{-1}(U_i) \subset \overline{\{x_j\}}$ pour
$j \geq i$ ; (Ep)$''$ $\overline{\{x\}}$ est une partie constructible de $X$,
c'est-à-dire l'image inverse de $\overline{\{x_i\}}$ pour $i$ grand ; (Ep)$'''$
$\overline{\{x\}}$ contient une partie constructible ouverte non vide.
L'encadré de la page 127 donne (Ep)$'$ $\Rightarrow$ (Ep), (Ep)$''$
$\Rightarrow$ (Ep)$'$ sous une condition de Mittag-Leffler et un ensemble
d'indices dénombrable, (Ep)$''$ $\Rightarrow$ (Ep)$'''$, et une réciproque
(Ep)$'''$ $\Rightarrow$ (Ep)$'$ marquée « ??? »\footnote{La première
implication est démontrée page 125 : $\varphi_{ij}^{-1}(U_i)$ est irréductible
de point générique $x_j$. Les autres ne le sont pas ; la flèche de la
deuxième porte « M.L.\ et $\Sigma = \mathbf{N}$ », celle de la troisième « 4 ». Le
reste de la page 127 est fait de calculs épars sur des idéaux premiers, sans
énoncé.}. L'idée est qu'un point épais est un point générique d'une partie
« de dimension finie » d'un objet de dimension infinie, comme le sont les
schémas de Greenberg d'un schéma sur un anneau de valuation discrète.

\subsection*{Applications rationnelles entre schémas de Greenberg (pages 130 et 131)}

Soit $V$ un anneau de valuation discrète de corps résiduel parfait $k$ et de
corps des fractions $K$. Pour $X$ sur $V$, $\overline{X}$ est le schéma de
Greenberg, avec $\overline{X}(k') = X(V_{k'})$ pour $k'$ algébrique sur $k$,
$V_{k'}$ l'extension correspondante de $V$, de valuation discrète si $k'$ est
parfait. Soit $\varphi : X_K \to Y_K$ un $K$-morphisme, $Y$ propre sur $V$. Il
définit une \emph{application d'ensembles} $\bar\varphi : \overline{X} \to
\overline{Y}$ : un point $\bar x$ donne, avec $k'$ la clôture parfaite de
$k(\bar x)$, un point de $X(V_{k'}) \subset X(K_{k'})$ ; son image par $\varphi$
est dans $Y(K_{k'}) = Y(V_{k'})$ (critère valuatif de propreté), c'est-à-dire
dans $\overline{Y}(k')$, et l'on prend le point sous-jacent. On a $k(\bar\varphi(\bar
x)) \to k(\bar x)^{\mathrm{parf}}$.

\textbf{Proposition 1}, \emph{telle qu'elle est énoncée.} Pour $X$, $Y$ intègres
de type fini, $\bar x$ épais et $\varphi$ dominante et génériquement séparable,
$\bar y = \bar\varphi(\bar x)$ est épais, $k(\bar y) \subset k(\bar x)$, et
$\bar\varphi$ provient, sur un voisinage constructible de $\bar x$, d'un morphisme
« d'indice fini ». \emph{Corollaire} : l'ensemble des points épais de
$\overline{X}$ est fonctoriel en $X_K$ pour les applications rationnelles
dominantes génériquement séparables entre $K$-schémas propres.
\textbf{Proposition 2} : si $X_K$ est lisse et $S$ une partie constructible de
$\overline{X}$, $\bar\varphi(S)$ est constructible, et $S$ admet une partition
finie sur les morceaux de laquelle $\bar\varphi$ est morphique d'indice fini ;
d'où la fonctorialité de l'ensemble $\mathrm{Cons}(\overline{X})$ des parties
constructibles.

Les démonstrations ne sont pas sur les pages, et « d'indice fini » n'est pas
défini ; il faut sans doute entendre : défini à un niveau fini de troncation de
Greenberg\footnote{Lecture mince : dans la proposition 1, « dominant
génér.t » et ses corrections interlinéaires (lues « excité », « en
$\bar x$ ») sont incertaines. Une marge propose, quand $Y$ n'est pas propre,
une « application non partout définie ».}. C'est ici que les applications
constructibles des pages 81 à 115 trouvent leur emploi : $\bar\varphi$ n'est
pas un morphisme, mais il l'est par morceaux. Le dossier 43 (pages 91 à 97)
construit de même, sur un trait, un foncteur « dit de Néron » par limite sur
les partitions constructibles\footnote{Le mémoire de Néron, « Modèles minimaux
des variétés abéliennes sur les corps locaux et globaux », \emph{Publ. Math.
IHÉS} 21 (1964), construit les modèles en travaillant sur les points à valeurs
dans les extensions non ramifiées par le foncteur de Greenberg ; nous ne
savons pas si le mot « épais » y figure, et ne l'affirmons pas.}.

\subsection*{L'ind-schéma $A \mapsto A((t))^{*}$ (pages 133 et 134)}

Le foncteur $\Gamma : A \mapsto A((t))^{*}$ des unités des séries de Laurent est
ind-représentable par $\mathcal{G} = \varinjlim_n \mathcal{G}_n$, où
$\mathcal{G}_n$ représente les couples de séries $\sum a_i t^i$, $\sum b_j t^j$
inverses l'une de l'autre avec $a_i = b_i = 0$ pour $i < -n$ : un sous-schéma
fermé de $\mathrm{Spec}\,\mathbf{Z}[(a_i, b_i)_{i \geq -n}]$. Pour $n = 0$,
$\mathcal{G}_0 = \mathbf{G}_m \times \mathrm{Spec}\,\mathbf{Z}[(a_i)_{i \geq 1}]$
représente $A[[t]]^{*}$. On a $\varphi_n : [-n, n] \times \mathcal{G}_0 \to
\mathcal{G}_n$, $(\nu, u) \mapsto t^{\nu}u$.

\textbf{Théorème} (page 134). \emph{(i) $\varphi_n$ est une immersion fermée
surjective. (ii) Sa restriction à $\{\nu\} \times \mathcal{G}_0$ est une
immersion ouverte pour $\nu = \pm n$, et ne l'est pas pour $|\nu| < n$, après
aucun changement de base non vide. (iii) $\mathcal{G}_n$ est la somme disjointe
d'ouverts fermés $\mathcal{G}_n^{\nu}$ de même espace que les images, et
$\mathcal{G}_n^{\nu} \to \mathcal{G}_{n'}^{\nu}$ n'est jamais un isomorphisme.}

(i) se vérifie : l'image de la composante $\nu$ est le fermé $a_i = 0$ ($i <
\nu$), $b_i = 0$ ($i < -\nu$), sur lequel $a_\nu b_{-\nu} = 1$, et ces fermés sont
disjoints ; elle est surjective parce qu'une unité de $k((t))$, $k$ un corps,
est un $t^{\nu}$ fois une unité de $k[[t]]$. Pour $\nu = n$, l'image est l'ouvert
où $b_{-n}$ est inversible (les coefficients de $t^{-2n}, t^{-2n+1}, \ldots$
dans le produit tuent alors les $a_i$, $i < n$) ; de même pour $\nu = -n$. Pour
$|\nu| < n$, les éléments $1 + \varepsilon t^{-1}$, $\varepsilon^2 = 0$, montrent
que l'image n'est pas ouverte. Le théorème dit, en termes d'aujourd'hui, que le
groupe des lacets de $\mathbf{G}_m$ a pour réduit $\mathbf{Z} \times
L^{+}\mathbf{G}_m$, et que la structure non réduite est portée par les
« queues » nilpotentes\footnote{C'est la description du groupe des lacets de
$\mathbf{G}_m$ qu'on trouve chez Contou-Carrère (\emph{C. R. Acad. Sci.} 1994) et
chez Beilinson--Drinfeld ; références citées de mémoire, rapprochement nôtre.
Une marge dit « $\varinjlim$ par abus de langage » : $\mathcal{G}$ n'est pas un
schéma.}.

\subsection*{Pro-schémas (page 135)}

Soit $\mathrm{Pro}''(\mathrm{Sch}/S)$ la catégorie des pro-schémas $(X_i)$ à
transitions affines à partir d'un rang, avec $X_i$ quasi-compact et de
présentation finie sur $S$ pour $i$ grand. Pour $\mathfrak{X}$ dans
$\mathrm{Pro}''$ et $\mathcal{Y} = (Y_j)$ à transitions affines, $Y_j$
quasi-compacts et quasi-séparés,
\[
  \mathrm{Hom}_{\mathrm{Pro}(\mathrm{Sch}/S)}(\mathcal{Y}, \mathfrak{X})
  \xrightarrow{\ \sim\ } \mathrm{Hom}_S\bigl(\varprojlim \mathcal{Y}, \varprojlim
  \mathfrak{X}\bigr) :
\]
le passage à la limite est pleinement fidèle sur $\mathrm{Pro}''$ (EGA IV
8.13.1)\footnote{La page ne demande aux $Y_j$ que d'être quasi-compacts ;
« quasi-séparés » est nécessaire à EGA IV 8.13.1 et on l'ajoute.}.

\emph{Remarque} : « Dans Serre et Néron, on utilise des systèmes projectifs de
schémas (de Greenberg) qui sont dans $\mathrm{Pro}''$. Il y a avantage à les
interpréter comme étant de véritables schémas sur $k$, au lieu de
pro-objets. » C'est la remarque que cite le résumé\footnote{Serre :
\emph{Groupes proalgébriques} (Publ. Math. IHÉS 7, 1960), cité de mémoire.}.

\subsection*{Groupes de changements de coordonnées (pages 137 à 141)}

Les dernières pages sont des calculs. Soit $A$ local d'idéal maximal
$\mathfrak{m}$, d'« échelon infinitésimal » $N$ ($\mathfrak{m}^{N+1} = 0$). Un
système $f = (f_1, \ldots, f_n)$ de séries
\[
  f_i = a_i + \sum_j b_{ij} t_j + \sum_{|p| \geq 2} c_{ip} t^{p}, \qquad a_i \in
  \mathfrak{m},\ \det(b_{ij}) \notin \mathfrak{m},
\]
définit une bijection $\lambda_f$ de $\mathfrak{m}^{n}$ sur lui-même, qui ne
dépend que des termes de degré $\leq N$. Les changements de coordonnées sans
terme constant forment le schéma en groupes
$G_{n,N} = \underline{\mathrm{Aut}}_{\text{alg.\ aug.}}\,\mathbf{Z}[t_1, \ldots,
t_n]/(t_1, \ldots, t_n)^{N+1}$, dont le transformé de Greenberg $\Gamma_n$ opère
sur $\mathfrak{m}^n$ regardé comme schéma sur $k$ ; avec les translations
$a \in \mathfrak{m}^n$ on obtient le groupe $\mathbb{G}_{A,n}$ des changements de
coordonnées affines-polynomiaux de la page 141. L'échelle $\mathfrak{m}^k /
\mathfrak{m}^{k+1}$ découpée en tranches $V_1, \ldots, V_N$ (page 137) et les
termes $\mathcal{T}^{0}_{B_0/A_0} \otimes V_i$ de la page 141 suggèrent un
dévissage de ce groupe par des groupes vectoriels, faisceau tangent tensorisé
par les gradués\footnote{Conjecture de lecture, sur une base très mince : la
lettre lue $\mathcal{T}$ est une conjecture du transcripteur, et le calcul de la
page 139 s'interrompt sans conclusion.}. La page 141 calcule enfin, modulo
$t^{3}$, la composition de deux polynômes $a'_0 + a'_1 u + a'_2 u^2$ et $u = a_0
+ a_1 t + a_2 t^2$ : les coefficients sont $a'_0 + a'_1 a_0 + a'_2 a_0^2$, $a'_1
a_1 + 2a'_2 a_0 a_1$ et $a'_1 a_2 + a'_2(a_1^2 + 2 a_0 a_2)$, ce qui est exact. Le
dossier s'arrête là.

\section*{Ce que le dossier annonce et ne fait pas}

\begin{itemize}
  \item Les deux « nous admettrons » des pages 19 et 20, qui laissent
  ouverte la démonstration du théorème de la page 16 sous sa forme générale.
  \item La démonstration du théorème de la page 25, que la lettre propose
  d'envoyer et que le dossier ne contient pas.
  \item La question de la page 38 (commutation de $f^{*}$ aux intersections
  filtrantes) et celle de la page 96 ($\mathrm{Cons}$ et limites projectives).
  \item La « catégorie géométrique » de la page 98, dont la page 100 ne donne
  que les axiomes.
  \item Les « espaces de Néron » du plan de la page 117, et la notion
  d'« indice fini » des pages 130 et 131.
  \item La rubrique 9 du plan de la page 33 (cas d'un $S' \to S$ étale) et les
  compléments de Greenberg de la rubrique 10.
  \item Le calcul de la page 139, interrompu.
\end{itemize}

\end{document}
